\section{Different intermediate histories can reach the same closing state}
\label{sec:native-conditional-sequence}

The preceding quantity restriction excludes a return of goods that were never
sold. A feasible monetary assignment can still describe a payment that the
company cannot fund, or a sale of a lot it does not hold. Applying the method
therefore requires a complete opening state and an ordered change to that
state. Keeping several possible carrying amounts also requires the later sale
to use the amount created by its own preceding settlement.

We retain the same synthetic return, now within one complete opening book.
Opening cash is 400 dollars and other inventory is 60 dollars. The refund
liability and recovery asset are 40 and 16 dollars respectively. All other
balances, component identities, existing lots, tax memory and prior event
identities are supplied explicitly in the retained opening. They are not
reconstructed from those four amounts. The two events are settlement on
1 February 2025 followed by sale of the newly recovered lot. The default case
orders both events on that date; another checked case sells on 1 March.

\subsection{Apply the conditional amount to a complete opening}

The \path{ConditionalReturnStateFamily/1} definition contains the complete
opening \path{state}, the selected \path{event}, numeric \path{event_domains},
a proposed \path{premise}, remaining \path{obligations} and a finite
\path{branch_limit}. A settlement domain may vary an explicitly unknown
returned quantity, unit handling cost, unit deterioration or recoverable value.
For a sale, only explicitly unknown proceeds may receive such a domain.
A domain cannot replace a known event fact. Each domain retains its exact
interval, premise and reference to the original event record and field.

Dates, identities, fractions, Boolean findings and externally supplied cash
facts keep their given known or unknown status. In particular, the solver
cannot decide that return rights have closed, invent performance of a sale,
or change refund cash to make its equations feasible. The original builders
regenerate the full event-dated monetary types, prerequisite bindings and
source references. An unknown required amount in the opening state remains
unresolved in this fixed-opening profile.

The compiler preserves every selected prerequisite and every case of the
carrying-value minimum. It also checks the native structural conditions:
the event identifier must be new, a settlement cannot overwrite an existing
lot, and a sale requires an existing undisposed lot. A valid opening does not
by itself guarantee that the next event is funded. Write \(C_0\) for opening
cash and retain the earlier refund and handling amounts \(R\) and \(H\).
The cash after settlement, denoted \(C_1\), must satisfy
\begin{equation}
\label{eq:native-family-funding}
 C_1=C_0-R-H\geq0.
\end{equation}
The equality is the cash posting identity; the inequality is the selected
native funding condition. In the example, \(400-40-8=352\) dollars remains.
If a valid opening instead contains only one dollar of cash, every proposed
settlement consistent with those cash facts fails the condition. The resulting
conditional family is empty. An opening that is already cash-negative has
the original \path{funding_failure} status before this family is constructed.
These are different questions, and neither creates an automatic loan.

\subsection{Carry the same lot cost through the sale}

The complete two-event definition, \path{SequentialReturnStateFamily/1},
contains the complete \path{first_family}, a fixed \path{sale_event}, an optional
\path{proceeds_domain}, a continuity \path{premise}, \path{observations} and
remaining \path{obligations}. Its first event must be settlement and its second
event must sell the whole newly created lot. The sale date cannot precede
settlement or fall outside the opening book's period; its identity must differ
from the settlement and all prior events. Unknown performance or fraction
remains unresolved. False performance and a fraction other than one retain
their unsupported consequence in this restricted sequence.

No intermediate state is selected to construct the second event. Instead,
every first-event branch retains its original variables and constraints.
The sale's cost is the same symbolic carrying amount \(V\) created in that
branch. The original sale expressions provide its cash, revenue and inventory
postings. This lowering supports the inspected references, constants, addition
and subtraction; a changed definition outside that syntax is unsupported.
The explicit continuity record identifies the original family, the carrying
calculation, the lot and both dates. Thus carrying amount at settlement becomes
cost of the same lot at sale under a stated correspondence, without silently
removing the date from either monetary type.

Let \(P\) be sale proceeds in dollars, let \(I_0=60\) be the inventory already
held, and use subscripts 1 and 2 for the states after settlement and sale.
Selling the entire new lot removes exactly the inventory added at settlement:
\begin{align}
\label{eq:native-sequence-balances}
 I_1&=I_0+V, & I_2&=I_1-V=I_0,\\
 C_1&=C_0-R-H, & C_2&=C_1+P.
\end{align}
Both cash balances must be nonnegative and proceeds must be nonnegative.
A large later receipt cannot repair a negative \(C_1\). Even when the events
share a calendar date, their explicit order preserves that intermediate test.

The refund and handling facts still fix \(n=2\). With \(N=20\), the first
minimum case gives \(V=20\); with \(N=5\), the second gives \(V=5\).
These are two representative assignments within the retained family. Total
inventory after settlement is therefore 80 or 65 dollars, not 20 or 5 dollars.
With \(P=40\), both assignments end with cash 392 and inventory 60 dollars.

Their profit paths differ too. Let \(\pi_1\) denote cumulative profit after
settlement from the supplied zero opening income and expense balances, before
any later tax-closing operation. The retained journal has revenue zero,
cost of goods sold \(16-V\) and selling expense 8 dollars. The sale then
adds revenue \(P\) and cost of goods sold \(V\). Consequently,
\begin{align}
\label{eq:native-sequence-profit}
 \pi_1&=-(16-V)-8=V-24,\\
 \pi_2&=(V-24)+(P-V)=P-24.
\end{align}
The carrying amount cancels from cumulative profit after the whole-lot sale.
At the two representative assignments, \(\pi_1\) is \(-4\) or \(-19\)
dollars, while \(\pi_2=16\) dollars in both. The retained replay also finds
the same complete terminal native state, including the disposed lot and tax
memory. The journals and intermediate states remain different. This
cancellation follows from the selected two-event treatment; it does not remove
the need to preserve the path or establish a general result about other methods.

\subsection{Use observations without replacing the family}

An observation contains its identifier, event prefix 1 or 2, named account,
complete type and an exact interval with source reference and premise. Its date,
entity, book, currency, dimensions and scale must match that prefix account.
Requiring total inventory to equal 65 dollars after settlement restricts
\(60+V=65\), hence \(V=5\). This excludes the first minimum case and retains
the second. Requiring inventory 60 after sale retains both cases because
equation~\eqref{eq:native-sequence-balances} already gives that closing value.

Replacing the first family by its first witness, with \(V=20\), would wrongly
discard the path consistent with the intermediate observation. The sequence
instead intersects the observation with every complete first branch. Every
new domain, observation and continuity premise remains in the result. Missing
premises leave the formulation incomplete and unresolved; remaining source or
state obligations keep actual history existence unknown.

For a feasible assignment, replay materializes only the declared unknown
numeric event fields. The repaired one-event family and both original native
transitions check the complete states, journals and flows. The original sale
monetary evaluator, every compiled prefix balance, observation and constraint
must agree exactly. An alternative valid assignment passes the same checks;
matching the solver's first witness is unnecessary. Saved JSON results replay
after serialization, including the original source bindings.

The global branch, cell and elapsed-time policies retain their earlier meaning.
One valid path establishes conditional existence, while unvisited branches
leave coverage partial. In the intermediate-inventory example, examining only
the rejected first case yields unknown, not an empty family, because the valid
second case remains unvisited. Emptiness requires checked infeasibility of every
retained case. Preparation and replay are bounded Python procedures rather than
hard interruption at the elapsed deadline.

The complete definitions, adverse cases and saved replays are under
\path{data/accounting-master/}:
\begin{quote}\small
\path{conditional-return-state-family-20261024-r2}\\
\path{sequential-return-state-family-20261024}
\end{quote}
The repaired first package preserves its failed serialization predecessor as
historical evidence. Both successful profiles describe hypothetical, uncommitted
states; source admission, accounting admission and method eligibility remain
false. Arbitrary event sequences, variable lot collections, source-based legal
and economic state, and wider accounting methods still require separate work.

\section{Ask what follows from the complete family}
\label{sec:family-outcome-queries}

The retained histories answer a financial question that a single simulated
path cannot settle: which amounts could occur, and which amounts must occur
under the supplied assumptions? In the preceding example, inventory of
65 dollars after settlement is possible. Inventory of 60 dollars after the
whole-lot sale is necessary. The distinction depends on all admissible paths,
including paths that a first witness does not select.

\subsection{Separate a question from an observation}

An observation of inventory of 65 dollars restricts the historical family. Asking whether
inventory of 65 dollars is possible leaves that family unchanged. The query therefore
carries no additional source assertion. It retains the original opening,
event order, branch conditions, funding restrictions, observations and premises.
It also retains unresolved obligations concerning the actual company.

Write \(\mathcal H\) for this complete conditional family and \(h\) for one
history within it. At a chosen prefix \(t\in\{1,2\}\), let
\(A_{a,t}(h)\) be the signed balance of account \(a\) after that event.
The query asks about the scalar amount
\begin{equation}
\label{eq:family-outcome-amount}
 f_t(h)=\sum_{a\in\mathcal A}c_a A_{a,t}(h),
 \qquad h\in\mathcal H.
\end{equation}
Here \(\mathcal A\) is the set of distinct queried accounts, and each
\(c_a\) is an exact rational, dimensionless coefficient. Every account must
have the same full monetary type at that prefix: entity, legal book, date,
currency, dimensions and scale. A threshold has that same monetary type.
This is a same-date account sum, not permission to add differently dated or
differently denominated amounts. The retained profit query applies the signed
account convention to revenue, cost of goods sold and selling expense.

\path{SequentialOutcomeQuery/1} contains the full \path{family}, the chosen
\path{prefix}, the complete \path{type}, the account \path{terms}, an
\path{operator} and a rational \path{threshold}, alongside its \path{profile}
identifier. The five comparisons are equality, strict less or greater, and
weak less or greater. Arbitrary callbacks and compound expressions are outside
this first interface.

\subsection{Distinguish possible, necessary and empty}

Let \(Q(h)\) denote the selected comparison of \(f_t(h)\) with its threshold.
The implementation asks two existence questions on the unchanged family:
\begin{equation}
\label{eq:family-outcome-partition}
 \mathcal H_Q=\{h\in\mathcal H:Q(h)\},\qquad
 \mathcal H_{\neg Q}=\{h\in\mathcal H:\neg Q(h)\}.
\end{equation}
A checked member of \(\mathcal H_Q\) establishes possibility. A checked
member of \(\mathcal H_{\neg Q}\) is a counterexample to necessity. If both
exist, the comparison is contingent. If the first exists and the second is
proved empty, it is necessary within this family. If the first is proved empty
and the second exists, it is impossible. When both sets are empty, the answer
is \path{empty_family}; necessity is not applicable. In particular, an
unfunded settlement does not make every financial assertion meaningfully true.

For inventory of 65 dollars after settlement, a path with \(V=5\) is a satisfying
witness and a path with \(V=20\) is a counterexample. For terminal inventory of 60 dollars,
equation~\eqref{eq:native-sequence-balances} supplies the invariant, while
the complete branch check proves that no counterexample meets the original
conditions. Inventory of 64 dollars after settlement is impossible in this example.
Adding the observation inventory of 65 dollars first would instead make that amount
necessary in the smaller, observed family. These are different questions with
different conditioning information.

Every witness replays both original native transitions. The verifier computes
the account sum independently with exact rational arithmetic and checks which
side of the comparison it satisfies. An empty region requires complete
infeasibility evidence and a separately checked unsatisfiable constraint core.
The checker accepts a different valid witness or a different valid core;
reproducing the solver's first choice is not the criterion.

Search is divided into elementary regions. For equality, the regions are
equal, less and greater; for each inequality, its complement uses the exact
opposite strict or weak comparison. One shared cell-count and elapsed-time budget spans all
regions of all parent branches. An unvisited region cannot prove emptiness.
A found counterexample can already refute necessity, even if the search is
otherwise incomplete; the result preserves that fact while leaving unproved
conclusions unknown. The number of witnesses supplies no probability.

\subsection{Bound the amount without selecting a history}

The closed \path{SequentialOutcomeRange/1} request contains its \path{profile},
\path{family}, \path{prefix}, \path{type} and account \path{terms}. It asks
for the tight lower and upper endpoints instead of a comparison. For a
nonempty family these are
\begin{equation}
\label{eq:family-outcome-envelope}
 L=\inf_{h\in\mathcal H} f_t(h),\qquad
 U=\sup_{h\in\mathcal H} f_t(h).
\end{equation}
An endpoint is reported as attained only if an original admissible history
reaches it. Thus an open interval has valid finite endpoints without supplying
executable states at those endpoints.

The supported fragment has finitely many affine branches with explicitly
finite numeric input domains. Each feasible branch is a convex region defined
by rational linear equalities and strict or weak inequalities. Within a
nonempty such branch, a closure point can be approached from an original
feasible point: their convex combinations satisfy every strict inequality
until the limiting closure point is reached. Optimizing over the closure can
therefore locate its infimum or supremum. Certification additionally checks
that the endpoint is feasible in that closure, no original value is better,
and the original strict constraints either admit or exclude endpoint equality.
An attained endpoint must also pass complete native replay.

The global lower endpoint is the smallest certified branch infimum; the upper
is the largest certified branch supremum. This aggregation requires every
original branch to be either checked infeasible or certified at both endpoints.
A partial search returns local branch evidence with no certified global bounds.
An empty family has no bounds. An unbounded answer is invalid in this explicitly
finite fragment, rather than an authorized infinity result.

The amounts in the example follow directly from the shared carrying value.
With returned quantity fixed at two, \(V=\min(20,N)\). The supplied recoverable
value domain is \(5\leq N\leq30\), so it permits both minimum branches and the carrying
envelope is \([5,20]\). Substituting into
equations~\eqref{eq:native-sequence-balances} and
\eqref{eq:native-sequence-profit} gives
\begin{align}
\label{eq:family-outcome-example}
 I_1&=60+V\in[65,80], & I_2&=60,\\
 \pi_1&=V-24\in[-19,-4], & \pi_2&=P-24.
\end{align}
At proceeds \(P=40\), cumulative terminal profit is 16 dollars before tax
closing. If proceeds instead range over \([0,100]\), its envelope is
\([-24,76]\). Opening inventory of 60 dollars is included in both inventory queries.
These equalities use the original signed postings; selecting two convenient
witnesses would not establish the bounds.

If recoverable value is restricted to \(5<N<20\), then \(V=N\), and the
intermediate inventory envelope becomes \((65,80)\): neither endpoint is
attained. In the original two-branch family, both branches support the upper
bound of 80 dollars, but only the branch with \(V=20\) attains it. The other approaches 80
from below. Accordingly, \path{supporting_branches} and
\path{attaining_branches} are different result fields. Global attainment
requires at least one attaining branch at the global endpoint.

Even a tight outer envelope need not describe every feasible amount. A union
of branches could permit only 65 and 80, leaving the intervening values
unattainable while retaining the same endpoints. The range interface therefore
reports an envelope and endpoint membership, not a general proof of interior
feasibility. A predicate query can separately test a particular interior amount.

The reference definitions and saved complete requests/results are in
\path{data/accounting-master/}:
\begin{quote}\small
\path{sequential-outcome-query-20261024}\\
\path{sequential-outcome-range-20261024}
\end{quote}
Both interfaces preserve missing premises and unresolved target obligations.
Where those obligations remain, an answer concerns the candidate family and
the actual outcome stays unknown. Their states are hypothetical and uncommitted;
source admission, accounting admission and method eligibility remain false.
General histories, nonlinear or integer-lot methods, actual-company state
reconstruction and independent implementations remain outside these checked
references.

\section{Retain the evidence needed to reproduce an answer}
\label{sec:simulator-evidence-archive}

An answer about a conditional family is useful only while its meaning can be
recovered. The reported inventory range alone does not identify the opening
book, permitted recoverable values, event order or unresolved source obligations.
Retaining the complete definition and result allows a later reader to inspect
those conditions and a compatible implementation to repeat the calculation.
Storage must also distinguish a complete saved answer from an interrupted write.

AccDSL's earlier P07 reference already provides physical SQLite storage,
immutable records, evidence views, event replay, retry handling and one schema
migration for its original company engine. Its run construction and replay
instantiate that older engine directly. A newer native state can be saved as a
generic record without thereby becoming a state that the older engine can
execute. The \path{SimulatorEvidenceArchive/1} reference uses the existing
immutable nodes and transactions to retain evidence under a narrower, explicit
meaning: the original bytes have been archived and their relationships checked.

\subsection{Preserve exact records and their relationships}

A complete range result contains both financial values and diagnostic metadata,
including elapsed time. P07's original numeric contract refuses JSON floating
numbers; its identity conventions also normalize Unicode text and treat certain
named lists as sets. Those conventions serve its original interface but cannot
promise exact retention of every newer JSON record. For example, two spellings
of an accented name can have different UTF-8 bytes despite the same normalized
text, and two differently ordered constraint lists can receive the same generic
identity.

The archive retains the original JSON text as bytes inside a small closed
envelope. Hex encoding preserves those bytes without applying text normalization
or numeric conversion to their contents. A byte count and SHA256 digest detect
changes relative to the stored record. The declared profile label is also kept
as exact UTF-8 bytes. Whitespace, number spelling, list order, finite timing
metadata and unresolved fields therefore survive retrieval. Different original
bytes remain different records, even when an analyst considers their numerical
content equivalent.

The closed \path{SimulatorEvidenceBundle/1} contains its \path{profile} and an
ordered \path{records} list. Every record has a local \path{name}, a descriptive
\path{role}, a declared \path{profile}, the original \path{json} text and
\path{dependencies} naming earlier records in that bundle. Roles distinguish a
definition, state, proposal, result or note; none confers authority to execute.
Names must be distinct, and forward, missing, repeated or cyclic dependencies
are refused. Declared dependency order is preserved even though the underlying
physical dependency edges are unordered.

JSON must contain an object, have no duplicate keys, use valid UTF-8 scalar
values and contain no nonfinite numbers. When the payload states a profile, it
must equal its declared profile. An unfamiliar profile can still be retained
as unverified evidence. Callers supply explicit maximum record count, individual
byte size and total byte size. These limits control this reference's accepted
input; they are not recommendations about model size or performance. Validation
does not discard an inconvenient field or round an unsupported amount to make
the record fit.

For the retained range example, one bundle contains the complete opening state,
the query definition depending on that opening, and the complete answer depending
on both. Another preserves the candidate family's unresolved actual-source
obligation. The receipt identifies every record in order. A compatible later
query implementation can read the original request and repeat its semantic
checks; the archive itself supplies no new financial interpretation.

\subsection{Publish a complete bundle and recover it}

Inserting several records independently can leave an incomplete group if the
process stops between insertions. The archive therefore inserts the record
nodes and their common receipt in one transaction. The receipt depends on every
record, and each record depends on its named predecessors and the archive's
version marker. An exact retry produces the same receipt. Changed bytes produce
a new identity, rather than replacing an earlier record under the same name.

Before commit, the archive rereads the complete receipt and its records. This
step matters because an existing content identity alone cannot prove that stored
bytes remain intact. Retrieval checks closed envelopes, identities, byte counts,
hashes, dependency edges, ordering and the original JSON conditions. A corrupted
existing record therefore cannot pass an exact retry merely because its former
identifier is still present.

A new archive is prepared and marked in a temporary file in the destination
directory. After that file is closed, a non-overwriting hard link publishes it.
An existing unrelated database is inspected read-only and refused before the
older storage implementation can alter it. An unknown schema version or
incompatible archive marker is also refused; this reference introduces no
automatic migration. A hard process exit can leave an unreferenced staging file,
which the archive does not sweep away automatically.

The retained interruption checks observe no destination before initial
publication and a complete marked archive afterwards. For a bundle transaction,
exit before commit leaves no new receipt; exit after commit leaves the complete
bundle available on reopen. These are local process-recovery results. Power
loss, concurrent writers, distributed storage and filesystems without the chosen
publication operation require separate evidence. The unchanged P07 migration
checks continue to concern its own schema and engine.

\subsection{Check bytes separately from financial meaning}

Consider a proposed settlement whose closing cash is increased by one dollar
and whose signed original-equity balance is reduced by one dollar. The modified
book still balances. Generic storage can retain it and pass its byte-integrity
checks, but replaying the original settlement does not produce that state. The
original native proposal checker rejects it. This example separates three
questions: whether the bytes were preserved, whether the stated transition
produces them, and whether the underlying company facts and accounting treatment
have been admitted.

The archive answers the first question. Its result has \path{archived=true};
retrieval reports \path{byte_integrity_verified=true} after its checks. It keeps
\path{semantic_verified}, \path{source_admission}, \path{accounting_admission},
\path{method_eligible}, \path{simulation_executed} and \path{native_committed}
false. These are the archive's own conclusions, regardless of stronger assertions
inside a retained payload. A stored digest is an integrity comparison, not proof
of authentic evidence or a correct financial method.

The reuse investigation and reference cases are retained under
\path{data/accounting-master/}:
\begin{quote}\small
\path{simulator-storage-reuse-audit-20261024}\\
\path{simulator-evidence-archive-20261024}
\end{quote}
The reference preserves complete evidence across its tested restart conditions.
Persisting an executable native event cursor, migrating arbitrary histories,
joining sector methods and admitting real-company state remain separate design
and conformance work.


\section{Keep lease rights distinct within one company book}
\label{sec:lease-shared-book}

Preserving a proposal makes its calculations available for inspection. To
reproduce them, the company model must also preserve the rights that generated
the entries. A total lease liability and a total right-of-use asset do not tell
us how much of either balance disappears when one right ends. Nor does a
balanced calculation for one lease establish that the company can pay it after
another lease has used the same cash.

Consider the retained synthetic example with two rights in one lease. Right A
has a right-of-use asset of 90 and a liability of 100; right B has an asset of
30 and a liability of 50. Amounts are exact USD quantities in this example.
The supplied findings support tracing the amounts to their individual rights.
Terminating B therefore removes its asset of 30 and liability of 50, giving a
gain of 20. Allocating the total asset of 120 in proportion to B's share of the
total liability would instead remove
\[
  120\frac{50}{100+50}=40,
  \qquad 50-40=10.
\]
That gain of 10 is wrong for the stated traced-right inputs. The two rights
have different asset-to-liability ratios. Their totals conceal that difference,
so the simulator retains each right's amounts through termination and
remeasurement.

This example uses the inspected \path{CorporateResiduals/1} methods and their
historical synthetic source scope. The owner remains
\path{Synthetic-heterogeneous-rights}, with the Legal book, base branch, USD
currency and \path{historical_EU_IFRS_conditional} basis. The permitted execution
window begins on 1 January 2024 and ends immediately before 1 January 2026.
Complete opening balances, lease membership and method findings are supplied
premises. They are not inferred from an actual issuer's reports or evidence of
current accounting compliance.

\subsection{Follow the rights through the complete lifecycle}

Both opening rights have a supplied annual rate of 10 percent. On 1 January
2024, the example terminates B and replaces A's payment schedule with one
payment of 121 at the annual period end. The original method discounts that
payment for one year. For A, the liability rises from 100 to 110, and the
remeasurement adds the same 10 to its asset:
\[
  L_A^{(1)}=\frac{121}{1.1}=110,
  \qquad R_A^{(1)}=90+(110-100)=100.
\]
Here $L_A$ is A's positive liability amount and $R_A$ its right-of-use asset;
the superscript identifies the completed step. B is now zero. Combining B's
termination with A's remeasurement reduces the total asset by 20 and the
positive liability by 40. With debit-positive signed accounts, the journal is
therefore $\Delta\mathrm{rou}=-20$,
$\Delta\mathrm{lease\_liability}=+40$ and
$\Delta\mathrm{other\_gains}=-20$. The entries sum to zero; the credit to
gains records the 20 profit.

A second supplied modification occurs on the same date. It removes a further
10 of A's asset and 11 of its liability, producing a gain of 1. The remaining
amounts before remeasurement are asset 90 and liability 99. Replacing the
remaining annual payment with 99 gives:
\[
  L_A^{(2)}=\frac{99}{1.1}=90,
  \qquad R_A^{(2)}=(100-10)+(90-(110-11))=81.
\]
The remeasurement reduces the remaining liability by 9 and makes the same
reduction to the asset. Across this second step, the total asset falls by 19,
the positive liability falls by 20, and the additional gain is 1. Tracing the
removed amounts and then remeasuring the remainder explains all three changes.

On 1 January 2025, the supplied findings confirm one elapsed year with no
intervening accrual or modification. The payment of 99 comprises interest of
$90\times0.1=9$ and principal of $99-9=90$. Cash falls by 99, finance expense
increases by 9, and the signed liability increases by 90 to zero. The retained
cash-flow policy assigns 9 to operating outflow and 90 to financing outflow.
Finally, a supplied systematic depreciation amount of 81 removes the remaining
asset and records expense of 81. This depreciation amount is an input supported
by the example's finding; the wrapper does not infer a useful life.

With opening cash 200, closing cash is $200-99=101$. Both lease balances are
zero, and cumulative pretax profit from these four events is
\[
  \Pi=20+1-9-81=-69.
\]
The gains, interest and depreciation explain the profit independently of the
cash payment. This is a complete lifecycle for the supplied finite example;
it does not supply an annual tax close or a complete financial report.

\subsection{Apply each lease to the same cash balance}

Now instantiate a second synthetic lease with the same supplied right-level
amounts and its own identity and premise. A company opening containing both
leases must reconcile to assets of 240 and positive liabilities of 300. Cash
and equity belong to that one company opening. Adding an isolated lease
calculator's opening cash or equity would capitalize the company again each
time a method runs.

For example, opening company cash of 400 followed by one payment of 99 gives
301. Adding the component fixture's opening cash of 200 would give 501, an
incorrect result. The shared-book proposal instead applies only the original
journal deltas. It updates the selected lease's rights, preserves the other
lease's memory, and leaves unrelated company accounts unchanged. Right A in
the first lease and right A in the second lease are distinct: their identities
are the pairs of lease identifier and local right identifier.

Funding must be checked in event order. Under the explicitly supplied
\path{no_overdraft} policy, company cash of 300 funds both payments of 99 and
leaves 102. Cash of 150 funds the first payment and leaves 51. The second
payment then fails: its liability remains 90, while the first lease's liability
remains settled. The successful earlier events survive; the failed event adds
no journal, cash movement or after-state. This is a funding failure under the
company premise. The original component calculator's ability to calculate a
balanced entry ending in negative cash does not provide that funding policy.

\subsection{Bind proposals to complete state and original methods}

The \path{SharedLeaseBook/1} record contains context, original source-scope
hash, execution period, current date, the full original account registry, all
lease components, collection-completeness and instantiation premises, funding
policy, prior event identities, gross cash movements and the conditional-state
flag. Every right retains its asset, liability, rate and last measurement date.
Canonical rational strings preserve exact amounts. Before and after each event,
the signed balances sum to zero, the ROU account equals all named rights'
assets, and the negative lease-liability account equals all their liabilities.
Unknown membership cannot establish those equalities merely because a supplied
subset balances.

\path{LeaseBookEvent/1} binds a new event identity, lease identity, date, source
scope and one of the original \path{lease.reduce}, \path{lease.pay} or
\path{lease.depreciate} methods. Parameters must satisfy that method's closed
schema and required findings. Dates are nondecreasing within the execution
window; a selected right's measurement date cannot lie after the event.
Payment retains the original exact annual-anniversary requirement. A missing
required date does not create a dated prefix, and no leap-day convention is
silently supplied.

The permitted account changes are explicit. Reduction affects ROU, lease
liability and other gains. Payment affects cash, finance expense and lease
liability. Depreciation affects ROU and depreciation expense. The wrapper checks
these effects, reconciles the complete rights collection and enforces the
company's cash policy. Unknown supported inputs remain unresolved; malformed
records, false findings and unsupported schedules retain their distinct original
refusal categories. No invented loan or balancing entry repairs a refusal.

A \path{LeaseBookProposal/1} binds the complete before-state, event and original
method hashes to the hypothetical after-state, journal and gross cash effects.
Applying it recomputes the original method and compares the entire proposal.
A forged cash increase offset by equity can preserve balance while failing
this comparison. Context, source, period, unrelated accounts and nonselected
leases must remain unchanged. The result remains hypothetical:
\path{committed}, \path{source_admission}, \path{accounting_admission} and
\path{method_eligible} are false. Application here does not publish a database
transaction; an evidence-archive receipt preserves bytes without replacing
financial recomputation.

The original-scope audit and shared-book reference, including complete input
and output examples, are retained under \path{data/accounting-master/}:
\begin{quote}\small
\path{lease-module-composition-audit-20261024-r2}\\
\path{lease-shared-book-20261024}
\end{quote}
The checked composition preserves heterogeneous rights and one cash balance
for these methods. General lease valuation, FX, tax and report integration,
arbitrary module imports, and reconstruction of actual contract rights remain
separate design dependencies.


\section{One lease total can conceal different rights}
\label{sec:lease-allocation-family}

The preceding lifecycle starts with a supplied asset of 30 for right B.
Suppose that amount is unknown and only the combined right-of-use asset of 120
is given. Even if the two rights and their liabilities of 100 and 50 are known,
the total asset does not identify its allocation. Assigning 20 to B leaves 100
for A; assigning 30 to B leaves 90 for A. Both allocations produce exactly the
same opening signed account balances. Terminating B gives a gain of 30 in the
first case and 20 in the second.

Let $x$ be B's opening right-of-use asset, in the same exact USD units as the
previous example. A's asset is then $120-x$. Nonnegative asset amounts imply
\[
  0\leq x\leq120,\qquad R_A=120-x,\qquad R_B=x.
\]
This is a family of supplied hypothetical openings. Exactly two identified
rights, their liability allocation, the fixed 10 percent rate and the original
method findings remain premises. Actual aggregate disclosures need not establish
those premises. The family assigns no probabilities to its allocations.

\subsection{Find where the supplied lifecycle can complete}

Keep the original four events, including the annual payments and depreciation
of 81. Only the unknown opening asset allocation and B's corresponding traced
removal vary. Terminating B removes liability 50 and asset $x$, so its gain is
$50-x$. A's payment of 121 gives liability 110 and adds 10 to its asset. After
the first event, total ROU is therefore $130-x$.

The second reduction removes asset 10 and liability 11 from A. Remeasuring its
remaining payment of 99 gives liability 90, which is 9 below the intermediate
liability of 99. Its asset becomes
\[
  R_A^{(2)}=(130-x)-10+(90-99)=111-x.
\]
The original method requires a nonnegative remaining asset. For $x>111$, this
second reduction is outside its supported treatment. That refusal does not make
the opening allocation impossible; it identifies a method limitation under the
supplied event. The first completed event remains available.

For $x\leq111$, the payment of 99 settles liability 90, records interest expense
9 and leaves company cash 101. The final event then proposes depreciation of
81. It can be applied only when that expense does not exceed the remaining
asset:
\[
  81\leq111-x\quad\Longleftrightarrow\quad x\leq30.
\]
Thus the complete supplied sequence works on $[0,30]$. For $30<x\leq111$, it
reaches payment but rejects the final depreciation. Replacing 81 with a smaller
amount would change the supplied scenario; the reference keeps the refusal and
its successful prefix instead.

The three regions partition the original allocation interval:
\begin{center}
\begin{tabular}{lll}
\toprule
Allocation of B's asset & Outcome & Completed events \\
\midrule
$[0,30]$ & Complete supplied lifecycle & Four \\
$(30,111]$ & Depreciation amount rejected & Three \\
$(111,120]$ & Second reduction unsupported & One \\
\bottomrule
\end{tabular}
\end{center}
The endpoints matter. At $x=30$, depreciation removes the remaining asset
exactly. At $x=111$, the second reduction leaves zero ROU, so it succeeds, but
the later depreciation of 81 does not. Each failed event leaves the last
successful state unchanged. A response containing only the first region would
lose valid declared openings and conceal why their subsequent paths stopped.

\subsection{Keep financial outcomes on the same allocation}

Within the complete region, depreciation leaves asset $30-x$. Pretax profit
combines the first gain $50-x$, the second gain 1, interest 9 and depreciation
81. Write $R_{\rm end}$ for final ROU and $\Pi$ for cumulative pretax profit:
\[
  R_{\rm end}=30-x,\qquad
  \Pi=(50-x)+1-9-81=-39-x,\qquad 0\leq x\leq30.
\]
The possible asset amounts range from 0 to 30 and profits from $-69$ to $-39$.
Both depend on the same opening allocation. Subtracting the two expressions
eliminates $x$:
\[
  \Pi-R_{\rm end}=(-39-x)-(30-x)=-69.
\]
For example, $x=20$ produces asset 10 and profit $-59$; $x=30$ produces asset
zero and profit $-69$. Asset 30 and profit $-69$ are individually feasible
marginal endpoints, but they cannot occur together. The first requires $x=0$;
the second requires $x=30$.

A hypothetical additional observation of final ROU equal to 10 would select
$x=20$ and profit $-59$, conditional on completion of the supplied lifecycle.
The original unconditioned family still contains its rejected and unsupported
regions. The query records the additional observation as a supplied premise;
it does not present that number as a recovered issuer fact.

The interrupted paths also retain dependent financial amounts. In the
rejected-depreciation region, the last successful state has ROU $111-x$, profit
$42-x$, cash 101 and zero lease liability. In the unsupported-reduction region,
the first event leaves ROU $130-x$, profit $50-x$, cash 200 and signed lease
liability $-110$. These are prefix states, not completed-year statements. Keeping
their dates and failure positions prevents a partially executed path from being
reported as a complete result.

\subsection{Represent the family without dropping interrupted paths}

\path{LeaseAllocationFamily/1} binds the original source lock, a finite
allocation interval within $[0,120]$, explicit endpoint inclusion, an allocation
premise, its source reference, a complete-membership finding and a
\path{max_cells} budget. Amounts are canonical rational strings. The reference
intersects the declared interval with all three regions. Empty intersections
contribute no cell; every nonempty intersection remains visible. Missing
membership or a missing allocation premise returns unresolved.

For each checked cell, the response retains its domain, original outcome
category, completed-prefix length, failure index and reason, affine financial
expressions and a complete native replay. The replay contains opening state,
events, proposals and the last successful state. It checks a representative
allocation against the unchanged original methods. The inspected affine
derivation establishes the within-cell relationship; a few replayed allocations
alone would not prove a bound over a continuous interval.

The cell budget limits those native replays. If it is exhausted, the remaining
domains and categories stay present as unchecked cells. The response reports
partial search, supplies no complete-family bounds and makes no assertion that
all allocations complete. Once every retained cell is checked, any final-outcome
bounds are explicitly conditional on complete execution. A family consisting
entirely of interrupted paths has no such bounds.

\path{LeaseAllocationQuery/1} requires the explicit condition
\path{complete_execution}. Its outcome constraints intersect one shared
allocation interval, preserving dependence between asset, profit, cash and
liability. The response retains the unconditioned family alongside the selected
interval and its conditional ranges. An impossible combination has an empty
selection and no bounds. A separate \path{max_replays} budget controls the
query witness: with zero budget, a nonempty candidate interval remains visible
but its selection and bounds are not reported as verified.

The source locks and full recomputation reject changed methods or forged saved
results. Source admission, accounting admission, method eligibility, actual
history identification and production commit remain false. The exact original
dates, rates, schedules and depreciation are restrictions of this reference;
unknown rates or liabilities, arbitrary membership and broader lease valuation
require additional methods and cannot be handled by renaming this profile.

The derivation and complete saved requests and responses are retained under
\path{data/accounting-master/}:
\begin{quote}\small
\path{lease-allocation-family-audit-20261024}\\
\path{lease-allocation-family-20261024}
\end{quote}
They show how aggregate-compatible openings can remain distinct through a
supported accounting calculation, and how unsuccessful paths remain part of
the answer. Connecting that family to actual disclosed rights, measurements
and authority is still a separate source-reconstruction task.


\section{What actual lease disclosures identify}
\label{sec:henkel-lease-measurements}

The preceding allocation example starts with named rights, supplied contract
findings and a hypothetical opening book. A company report usually supplies
less. Before using it to select an opening, the simulator must distinguish what
was reported from the contract state needed to calculate the next payment or
modification. Henkel's retained 2025 annual report makes that distinction
concrete. It gives several useful lease measurements, but those measurements
neither enumerate the individual rights nor allocate their carrying amounts.

The following amounts are in millions of euros for the consolidated Henkel
Group. The report describes its basis as IFRS Accounting Standards as adopted
in the European Union. Its 2024 comparative and 2025 current figures belong to
the same retained 2025 document; their appearance together supplies no evidence
that these exact bytes were public in 2024. The basis note also says that
individual figures are rounded and addition can differ from the indicated
total. It supplies no rounding tie rule.\footnote{Henkel, \emph{Annual Report
2025}, p.~351. Page references in this section are printed numbers; physical
PDF pages are two higher. The source audit cited at the end binds the retained
PDF, its hash and exact source locations.}

At 31 December 2025, Henkel reports right-of-use assets of 594, lease
liabilities of 666 and undiscounted contractual lease cash flows of 721.
These amounts answer different questions. The first measures the remaining
carrying value of the rights to use assets; the second measures the recognised
liability; the third groups future undiscounted payments by maturity. The
liability consists of 557 noncurrent and 109 current. The maturity bands are
114 due within one year, 351 due between one and five years, and 255 due after
five years. Their displayed sum is 720, one below the reported total of
721.\footnote{Henkel, \emph{Annual Report 2025}, pp.~370--374, 397 and 423--424.
The maturity table is explicitly based on undiscounted cash flows.}

Different measurement bases explain why 594, 666 and 721 should not be forced
to agree. Rounding creates a separate issue even within a single basis.
The reported 2025 net right-of-use amounts for land and buildings, plant and
machinery, and factory and office equipment are 530, 16 and 47. Those three
numerical entries sum to 593; the total is 594. The construction category is
printed as a dash. The inventory keeps that dash distinct from zero and records
the difference between the total and the three numerical entries without
inventing a fourth amount. Similarly, the 2024 gross right-of-use amount is
1,173 and accumulated depreciation and impairment is 543, while net carrying
amount is 631. Treating every displayed integer as an exact underlying amount
would demand both $1{,}173-543=630$ and a net amount of 631. That exact-value
interpretation is wrong; it does not establish an error in the rounded report.

\subsection{Read the liability bridge with its signs intact}

The cash-flow reconciliation helps distinguish a payment from the expense
recognised in the same year. Henkel presents liabilities with a negative sign
in this reconciliation. A payment reducing the liability therefore enters
positively; an addition or accrued interest increasing it enters negatively.
For 2025, the displayed opening liability is $-713$. Cash changes are $158$,
comprising interest paid of 25 and redemption of 133. Interest expense is a
separate movement of $-25$, additions are $-139$, acquisitions or disposals of
subsidiaries contribute 1, foreign exchange contributes 50, and sundry changes
contribute 2. The fair-value movement is printed as \texttt{-0}.
Using the displayed numerical value of that last entry gives
\begin{equation}
\label{eq:henkel-displayed-liability-bridge}
 -713+158-25-139+1+50+0+2=-666.
\end{equation}
This is exact arithmetic on the displayed entries, in EUR millions, for the
calendar year 2025. The right side agrees with the negative of the reported
closing liability. The sign of the printed zero remains in the source record
even though its displayed rational value is zero.\footnote{Henkel,
\emph{Annual Report 2025}, pp.~450--451. The cash-flow statement on p.~347
reports redemption as an outflow of $-133$ in 2025 and $-146$ in 2024.}

The subtotal 158 enters equation~\eqref{eq:henkel-displayed-liability-bridge}
once. Adding 25 and 133 again would count the same cash twice. The interest
expense of $-25$ has a different role: it increases the liability that the
interest payment reduces. The 2024 reconciliation follows the same sign
convention, but the sum of its displayed opening and movements is $-712$,
whereas its displayed closing balance is $-713$. The retained residual is
$-1$. The rounding disclosure makes this a reason to examine precision, not
authority to insert a missing transaction.

Nor does the bridge explain every lease-related cash flow. Henkel reports
total cash paid for leases of 174 in 2025, exceeding the recognised-liability
cash subtotal of 158 by 16. Short-term lease expense, low-value lease expense
and excluded variable-payment expense are 11, 5 and 1, respectively. Their
sum is 17, and the report does not identify them as the allocation of that cash
difference. Expenses and cash flows can differ in both timing and scope. The
source inventory retains those measurements separately.

The same care is necessary for the asset movement. Opening net right-of-use
assets of 631, additions of 139, and depreciation and impairment of 132 give
$631+139-132=638$, compared with a reported closing amount of 594. Those three
movements are not a complete asset reconciliation. The difference of $-44$
does not identify a particular disposal, exchange effect or allocation to a
right. It remains a diagnostic of the incomplete comparison.

\subsection{Ask a conditional question about unrounded measurements}

The useful next question is whether one set of unrounded measurements could
satisfy selected relations while producing the displayed figures. Answering
it requires an explicit rounding hypothesis and explicit premises about which
underlying relations are exact. The report's general rounding statement does
not choose these premises for us.

For illustration, suppose that each selected amount is rounded to the nearest
million euros and exact half ties are rounded away from zero. Let $d$ denote
a displayed integer and $x$ the corresponding unrounded amount, both measured
in EUR millions. Under this proposed rule the inverse interval is
\begin{equation}
\label{eq:henkel-proposed-rounding}
 x\in I(d)=
 \begin{cases}
 [d-\tfrac12,d+\tfrac12),&d>0,\\
 (d-\tfrac12,d+\tfrac12],&d<0,\\
 (-\tfrac12,\tfrac12),&d=0.
 \end{cases}
\end{equation}
At the positive lower endpoint, a half tie rounds up to $d$; at the positive
upper endpoint it rounds to $d+1$. For a negative number the included endpoint
is reversed. Neither half endpoint rounds to zero. Thus the zero case is open
at both ends. Treating the printed \texttt{-0} numerically as zero imposes no
additional sign restriction: the source has not established what sign rule,
if any, its printed zero encodes.

This hypothesis can reconcile the 2024 gross, accumulated and net amounts
without changing their displayed values. For example, underlying amounts
1,173.25, 542.50 and 630.75 satisfy
$1{,}173.25-542.50=630.75$ and round to 1,173, 543 and 631 under the proposed
rule. The example is a conditional mathematical assignment, not a recovery
of Henkel's actual amounts. A different rounding convention would require
different endpoint rules.

Compatibility must also hold jointly. Let $L_{25}$ be the positive closing
liability, $L^{\mathrm{nc}}_{25}$ and $L^{\mathrm{cur}}_{25}$ its noncurrent and
current portions, $B_{25}$ its signed reconciliation closing, and $O_{25}$
the opening of the 2025 reconciliation. Let $B_{24}$
denote the signed closing at 31 December 2024. Selected underlying premises
include
\begin{equation}
\label{eq:henkel-shared-measurements}
 L_{25}=L^{\mathrm{nc}}_{25}+L^{\mathrm{cur}}_{25},
 \qquad B_{25}=-L_{25},
 \qquad O_{25}=B_{24}.
\end{equation}
The same $B_{25}$ appears as the closing amount in the 2025 stock-flow
relation, and the same $O_{25}$ appears there as its opening. Separate checks
could choose incompatible values for these shared amounts. One joint system
requires a common assignment. The opening is dated 31 December 2024, the
movements cover $[2025\text{-}01\text{-}01,2026\text{-}01\text{-}01)$, and
the closing is dated 31 December 2025; the opening symbol is indexed by the
reconciliation year, not its measurement year.

The checked reference combines nineteen such relations over fifty-four
selected measurements from the two years. With the proposed rounding and
identity premises, it finds an assignment satisfying every interval and
equation, checked with exact rational arithmetic. This establishes conditional
consistency of the selected measurements. It neither chooses the actual
unrounded amounts nor makes the rounded intervals independent. The 2024
displayed-as-exact counterexample fails because its three exact values cannot
satisfy the gross-to-net relation; that failed hypothesis is distinct from a
contradiction in the underlying accounts.

\subsection{Preserve the disclosures before admitting a lease book}

\path{HenkelLeaseMeasurements/1} binds this finite example to its retained
source. Its closed request contains \path{profile}, \path{inventory_sha256},
\path{selected_relations}, \path{rounding_hypothesis} and
\path{identity_hypotheses}. The hash selects the exact inventory. Relation IDs
must be unique members of the supported registry; the identity map contains
exactly those IDs. Each premise is absent or is the explicitly proposed
premise for that relation. Callers cannot replace the amounts, units, dates,
company or equations through this interface.

For each presentation year, the registry supports current plus noncurrent
liability; gross less accumulated right-of-use amounts; depreciation categories;
undiscounted maturity bands; cash subtotal components; the signed liability
stock-flow; the closing-liability sign comparison; the redemption sign
comparison; and the maturity table's separate carrying-value comparison.
Opening-to-closing continuity between the two years supplies the nineteenth
relation. Category sums involving unquantified dashes, the incomplete asset
movement, attribution of the cash difference to expenses, and equality between
discounted and undiscounted amounts are excluded.

Every response retains all 104 numerical-table and narrative entries, including
eight dash cells, the twelve qualitative observations and fourteen unresolved
state dependencies. The selected measurements carry their original source
references, literals, dates, units and consolidation context. A stock-flow
relation retains the roles of opening, period movements and closing; these
roles are validated before lowering to the existing rational constraint engine.
Shared observations keep one identity throughout the selected system.

By default, the adapter supplies no numeric inverse intervals. Missing rounding
or identity premises leave the assessment \path{unresolved}. An empty relation
selection reports \path{no_asserted_system} while retaining the source record.
A verified assignment gives \path{conditionally_feasible}; a checked
contradiction concerns only the declared conditional system. Resource exhaustion
returns \path{unknown}, and invalid source bindings or declarations return
\path{invalid}. The reference permits at most one conjunction cell, with an
explicit bounded solver budget; zero cell budget supplies no verified witness.
Saved results are rebuilt from the original request. Fixed fields must agree,
and saved witnesses must satisfy the exact intervals and equations; a different
valid witness or elapsed runtime need not have identical bytes.

An actual lease book still needs individual lease and right identities, complete
membership, legal lessees and owner books, unrounded allocations, contract and
functional currencies, rates, commencement and measurement dates, payment
schedules, modifications, depreciation and applicable findings. Henkel's
description of country-, duration- and credit-specific discount rates does not
supply those numeric inputs. Its description of optional lease periods does
not supply contract-specific exercise findings, and its qualitative
mid-triple-digit-million-euro payment description remains qualitative.
Short-term and low-value recognition exemptions require separate treatment.
Initial right-of-use measurement can also include payments, incentives and
other adjustments beyond the liability; equality between the two is not a
general allocation rule.\footnote{Henkel, \emph{Annual Report 2025},
pp.~373--374. The retained disclosure describes these policies; the present
example does not independently certify their application to each contract.}

The source adapter therefore leaves executable opening state absent and source,
accounting and native-state admission false. It preserves useful measurements
while keeping their missing contract explanation visible. The synthetic USD
lease methods retain their original company, period and authority scope; they
cannot become Henkel methods by changing a name or currency label. The full
requests, relations and exact checked results are retained under
\path{data/accounting-master/henkel-lease-measurements-20261024}, with the
reference in \path{scripts/henkel_lease_measurements_v1.py}. The source audit,
including the retained PDF path, SHA-256 and exact text and table locations,
is in \path{data/accounting-master/henkel-lease-source-audit-20261024}. Connecting these
measurements to actual contracts remains the next substantive reconstruction
problem.


\section{Return and lease methods in one company book}
\label{sec:shared-return-lease-book}

A customer refund and a lease payment draw on the same company's cash. Separate
method examples can each balance and still fail to describe that competition.
The lease example might spend money that the return example has already paid
to customers. Combining their closing reports would discover the inconsistency
too late: the second payment should have failed when it was proposed.
This makes a common book, rather than a common calling convention, the next
necessary part of the implementation.

The actual Henkel disclosures in the preceding section still lack the contracts
needed for an executable opening. The present example therefore declares a
new synthetic company. It uses retained return and lease calculations to test
shared cash, journals and memory, while leaving actual-company reconstruction
as a separate problem.

\subsection{A common caller preserves different capabilities}

The finite \path{ModuleInvocation/1} interface supplies one request form for
four existing references. Typed returns and shared leases propose transitions;
Henkel measurements assess conditional constraints; the sector reference
produces a report. Their outputs have different meanings. A measurement result
cannot become a posting merely because the caller uses the same envelope.

Each request names a registered module and operation, binds the complete
payload, and identifies the inspected implementation and schema bundle.
The caller selects code from its fixed registry. A supplied path or callback
cannot select another function. Typed-return inputs also retain their original
raw-byte hashes and duplicate-key checks. The complete native result survives,
including journals, gross flows, memory changes and refusal reasons.

An outer status of \path{returned} means that the selected module answered.
Its separate \path{native_status} may be proposed, unsupported, invalid or
unresolved. A request to post the sector report, post a measurement assessment,
or commit a lease proposal is refused before native invocation. Verification
recomputes complete transition and report results. The measurement verifier
instead permits another valid solver witness and another runtime while checking
the underlying equations and all fixed information.

Identity binding does not authenticate source labels. In the retained return
reference, changing a well-formed owner label or packet digest can still produce
a proposal; source admission remains false. The lease reference fixes its
original scope and refuses a changed scope digest. The common caller preserves
both behaviors. It cannot manufacture the missing source-to-state judgment.

\subsection{The original states are incompatible}

The return state belongs to synthetic owner P and begins in 2025. The lease
state belongs to a different synthetic owner and begins in 2024. Their schemas
are closed, their legal-book labels differ, and their account registries contain
22 and 26 accounts respectively. The return state stores recovered lots, tax
adjustments and fiscal history. The lease state stores named rights, their
allocations and a premise that all leases are listed.

Even shared account names are insufficient. Cash needs one owner, currency,
date and opening amount; inventory needs a complete allocation to stock; lease
liabilities need their rights and payment history. A field union does not
resolve those meanings and violates both original validators. Renaming the
lease owner does not repair the original scope restriction. The read-only
compatibility reference consequently rejects direct combination, preserves
unknown membership as unresolved, and accepts only its explicit control:
two valid identical snapshots of the same native profile. That control authorizes
no joint transition.

\subsection{One declared opening and explicit memory ownership}

\path{SyntheticReturnLeaseBook/1} supplies a separate synthetic USD legal book.
It retains the 22 return accounts and explicitly adds the right-of-use asset,
lease liability and lease derecognition gains. Debit balances and expenses are
positive; credit balances and income are negative. The six selected lease
posting categories---cash, right-of-use asset, lease liability, finance expense,
depreciation and other gains---have declared meanings in this 25-account
registry. No mapping between unrelated return cost categories is inferred.

The opening supplies cash 500, inventory 60, a recovery asset of 16, a refund
liability of 40, a right-of-use asset of 120 and a lease liability of 150, together
with the surrounding accounts recorded in the complete fixture. Its signed
original-equity balance is explicitly 20. These are synthetic premises chosen
for the example; the program neither infers that opening from an issuer nor
computes an equity plug. The two lease rights have right-of-use amounts 90 and
30 and liabilities 100 and 50. The declared collection contains exactly those
rights.

The book owns one period, date, cash balance, event history and gross-flow
history. Return methods own the return component, recovered lots and settlement
tax adjustments. Lease methods own the lease-right collection. Other memory,
including tax premises and prior fiscal history, must survive each selected
event unchanged. Exact write sets enforce that separation.

For a dated event $k$, let $B_k$ be the vector of signed account balances before
the event and $J_k$ its journal in the same account order, both measured in USD.
Let $F_k$ contain its gross cash flows, with receipts positive and payments
negative. Pure application checks the exact identities
\[
 B_{k+1}=B_k+J_k,\qquad
 \sum_a J_{k,a}=0,\qquad
 J_{k,\mathrm{cash}}=\sum_{f\in F_k}f.
\]
The first identity applies the journal once. The second preserves the signed
trial balance. The third prevents classified cash flows from disagreeing with
cash. Component allocations are checked as well: recovered lots explain their
inventory, and lease rights explain the right-of-use asset and lease liability.
An aggregate balance alone cannot replace those allocations.

Return settlement and whole-lot sale reuse the checked typed prerequisite,
monetary and memory programs. Their memory schema is extended explicitly to
retain all 25 balances and the lease fields. The lease operations reuse the
original component calculation kernel under the new declared context. The
fixed-owner lease wrapper remains unchanged. This is a separately specified
synthetic instantiation, not a lossless migration of the old owners.

\subsection{Follow both methods through the same cash balance}

At 1 January 2024 the lease reduction terminates right B and remeasures right A.
The journal reduces the right-of-use asset by 20, reduces the credit liability
by 40 and recognizes a gain of 20. No cash moves. The remaining right-of-use
asset is 100 and its liability 110.

On 1 February the return method pays customers 40 and handling costs 8,
releases the refund liability of 40 and recovery asset of 16, and recognizes
recovered inventory of 20. The associated cost-of-sales credit is 4 and selling
expense is 8. The two cash payments remain separate operating flows. A recovered
lot records the units, carrying amount, tax base and acquisition date, while
the settlement tax adjustment remains available for a later tax calculation.

On 1 March the entire recovered lot is sold for 30. Its carrying amount of 20
leaves inventory and enters cost of sales. On 1 January 2025 the lease payment
of 121 comprises interest 11 and principal 110 under the supplied annual
timing assumptions. Interest is an operating outflow and principal a financing
outflow under the retained election. A separately supported depreciation
amount of 10 then reduces the right-of-use asset. Neither that amount nor the
discount rate is inferred from aggregate disclosures.

With no other cash events in the example, the closing amount follows directly:
\[
 \mathrm{cash}_{\mathrm{close}}
 =500-\underbrace{(40+8)}_{\text{refund and handling}}
       +\underbrace{30}_{\text{recovered-stock sale}}
       -\underbrace{121}_{\text{lease payment}}
 =361.
\]
The book closes this sequence with right-of-use assets of 90, no remaining
lease liability and inventory of 60. Five event identities and five gross
flows survive. Return events leave the lease memory unchanged; lease events
preserve the return component, lots, tax adjustments and other return memory.

The low-cash case distinguishes genuine shared execution from two independent
examples. Start with cash 130 and omit the recovered-stock sale. The signed original-equity
amount is supplied as 390 in this alternative opening, keeping the declared
trial balance balanced; the program makes no automatic adjustment. The return
settlement leaves 82, so the later lease payment of 121 fails the shared funding
check. The completed lease reduction and return settlement remain; the failed
payment changes nothing. With the same opening cash but without the return
settlement, the lease payment succeeds and leaves 9. Cash competition therefore
changes the admissible sequence even though each method's arithmetic is
unchanged.

\subsection{What replay establishes}

A proposal binds the whole before-state, event and implementation bundle.
Application recomputes the proposal before returning a new in-memory state.
An old proposal cannot apply to a changed cash balance, and a forged journal,
flow, component change or admission flag fails full comparison. Event identities
cannot repeat. An empty event list still requires a valid opening. Unsupported
currency or tax operations are refused; missing return measurements or lease
membership remain unresolved.

All five selected event kinds are also compared with the original native
methods. The original return comparison stays in its 2025 period and original
owner; the new book expressly uses different return dates and ownership.
Journals and activity/amount cash flows match exactly for the same numerical
component inputs. Component and lot memory match apart from the declared lot
acquisition dates. The comparison supports reuse of those effects, not identity
between the old and new companies.

This reference supplies a finite shared-book execution example. It preserves
tax memory but does not close tax or produce complete financial statements for
the new profile. General module substitution, arbitrary histories, group and
foreign-currency methods, actual-source admission and physical database
durability remain separate requirements. The existing database reference is
tied to another engine; its transaction evidence does not transfer automatically.
The full schemas, complete openings, requests, proposals and checked comparisons
are retained under \path{data/accounting-master/} in
\path{module-invocation-20261024},
\path{shared-state-compatibility-20261024} and
\path{shared-book-successor-20261024}.

\subsection{Sell part of a recovered lot without losing its remaining units}
The whole-lot sale above leaves no recovered stock to track. A company that
sells only part of that stock needs a quantity and cost record for what remains.
The existing lot already records both; its carrying amount alone would not
identify how many units were sold. A separately selected method therefore
takes the sold quantity as an explicit input rather than deriving it from cash
proceeds or inventory value.

This method assumes that the remaining units are homogeneous and divisible,
that their measurement basis has not changed, and that proportional allocation
is the selected treatment. The caller also supplies affirmative findings for
completed performance and the retained equality of tax base and carrying
amount. Those findings record explicit premises. Unknown findings remain
unresolved, and false findings place the event outside this method's scope.
A heterogeneous lot or a different tax treatment requires another method.

Immediately before a sale, let $N>0$ be the remaining units, $C\geq0$ their
carrying amount in USD, and $T=C$ their tax base under the selected treatment.
The event supplies sold units $q$, with $0<q\leq N$, and cash proceeds $P\geq0$
in USD. The fraction applies to the stock immediately before this event:
\[
 f=\frac{q}{N},\qquad
 C_{\mathrm{released}}=fC,\qquad
 T_{\mathrm{released}}=fT.
\]
The remaining lot keeps its identity and acquisition date, with units $N-q$,
carrying amount $C-fC$ and tax base $T-fT$. It becomes disposed exactly when
$q=N$. A positive-quantity lot whose carrying amount is zero remains live;
zero cost does not imply that no goods remain.

The journal debits cash by $P$ and cost of sales by $fC$, and credits revenue
by $P$ and inventory by $fC$. Its debit-positive sum is
$P+fC-P-fC=0$. Cash rises by $P$ and inventory falls by $fC$, so net assets and
profit both rise by $P-fC$. In the indirect operating cash calculation, adding
back the inventory reduction to that profit effect gives
$(P-fC)+fC=P$, the same receipt as the direct cash record. The reduction in
the lot's carrying amount equals the inventory posting; every other lot and
the separate inventory balance remain unchanged. This sale assesses no tax
expense and makes no new measurement of inventory.

After the February settlement in the shared-book example, the recovered lot
has two units with carrying amount and tax base of 20. Cash is 452 and total
inventory is 80. Selling one unit for 15 releases 10 of cost, leaves one unit
with carrying amount and tax base of 10, and gives cash 467 and inventory 70.
The prior lease gain of 20, cost-of-sales credit of 4 and handling expense of 8
give recorded profit of 16 before this sale; its margin of 5 raises profit to
21. Operating cash since the opening is $-48+15=-33$. Assets of 827 less
liabilities of 214 equal equity of 613, also opening equity of 592 plus profit
of 21.

If the remaining unit is then sold for another 15 under the same conditions,
the new fraction is one: one unit is sold from one remaining unit. The second
sale releases the remaining 10, leaves cash 482 and total inventory 60, and
marks the lot disposed. The combined account changes equal those of the
original whole-lot sale for 30. The two event identities and cash receipts
remain separate; equal aggregate effects do not make their histories identical.

A proposal retains the entire before-state, event, allocation and resulting
book. Application recomputes the complete proposal. A changed remaining
quantity, altered allocation or balanced but different journal fails that
comparison. Selling more than the available units, selling a disposed lot,
reusing an event identity or backdating a sale is refused without changing
the successful preceding state.

This optional transition has checked quantity, journal and statement identities
and reproduces the original full-sale state when all units are sold. The
existing whole-lot method remains unchanged. Connecting the optional event to
the mixed-book authoring and annual dispatch interfaces is separate work;
these local controls do not establish that those interfaces already support it.

\subsection{Carry unsold stock through the assessed annual book}
The next composition supplies the annual execution step left open above. It
keeps the full tax-adjusted company book, including income tax payable,
assessment records, liability vintages, losses and earlier periods. The
standalone inventory method accepts 25 accounts, whereas this financial book
has 26. Removing tax payable or moving it into equity to reuse that smaller
schema would change the company. The annual extension instead validates and
updates the complete registered state.

The opening for this example comes from the checked retained-stock history,
with its planned whole-lot sale omitted. At 1 January 2027 it has two unsold
recovered units with carrying amount and tax base of 20, cash of 452 and tax
payable of $41/4$. New-period income starts at zero. Selling one unit for 15
on 2 January gives cash 467, inventory 70, a remaining unit with cost 10 and
recorded profit 5. Operating cash is 15. The payable remains $41/4$: an
inventory sale does not pay or reassess a prior tax liability.

The sale retains the complete assessment and liability histories. If it occurs
after a current assessment, that assessment's amounts remain recorded but its
status records subsequent activity. Any prior deferred measurement becomes
stale because it was bound to the earlier financial state. A sale inside the
already assessed interval is refused under the existing timing rule. New
activity is not silently inserted into a historical assessment.

Consider a hypothetical current assessment of 2 for the period containing the
first sale. It adds a separate liability, so total payable becomes $49/4$ and
recorded profit becomes 3. With explicitly supplied fresh deferred measurement
and the remaining carry conditions, a provisional boundary transfers that
profit to retained earnings and opens the next period with the unsold unit,
its cost of 10, cash 467 and the same liability vintages. Selling that unit for
15 on 2 January 2028 then gives cash 482 and recorded profit 5. That period's
operating cash is 15; the earlier receipt remains in the archived period.
The new-period payable is still $49/4$ until a payment or assessment changes it.

A separate sequence checks why fresh assessment matters. After a midyear
assessment and deferred measurement, an inventory sale records further
activity and invalidates the earlier measurement. Carry is unresolved at that
point. A fresh cumulative current assessment followed by a measurement bound
to the changed state permits the selected provisional carry; both earlier and
replacement assessments remain available. Reusing the old measurement after
that reassessment is invalid because its state digest no longer matches.
Attempting measurement before the required current reassessment instead stops
at that missing prerequisite. The two refusals answer different questions.

Ordinary operations continue through the protected annual evaluator. Their
financial states, reports and traces are unchanged; the new program profile
and request identities identify the extension. A full-quantity sale also
matches the old annual whole-lot sale on the complete resulting book, journal,
flows and report. Altering tax payable, an archived quantity or an allocation
fails complete-result verification even if the alteration is balanced.

The optional operation now has a checked annual execution meaning. Its matched
source declarations and annual-family dispatch still require separate
extensions. The current authoring and family profiles do not acquire support
merely because the financial transition exists. The examples remain conditional
on the stated inventory and tax assumptions; actual source and accounting
admission are unresolved.


\subsection*{Author inventory alternatives on the same annual book}

The optional annual inventory method can now be described in source text and
compared across named scenarios.  These interfaces extend the checked execution
above.  They retain its opening balances, tax liabilities and recognition
assumptions, so a change in how a sale is written cannot silently change the
company being simulated.

\texttt{Annual\allowbreak{}Inventory\allowbreak{}Authoring/1} uses the annual declaration grammar with an
explicit \texttt{partial\_\allowbreak{}inventory\_\allowbreak{}sale} operation.  Its header names the sale;
the body supplies the sale profile, date, allocation method, lot, units sold and
cash proceeds.  Each of the six inventory findings remains a separate
known-or-unknown declaration.  A known finding carries the text explaining its
provenance; an unknown finding carries a reason.  This inventory convention is
preserved alongside the structured provenance used by the older tax
operations.  Neither form is silently translated into the other.

Compilation binds the declared opening request to its exact content digest and
returns the same \texttt{Annual\allowbreak{}Inventory\allowbreak{}Program/1} request as the typed form.
It does not execute a sale, assess tax or advance a period.  Consequently, a
false divisibility finding or an unresolved homogeneity finding survives as
explicit data until the financial evaluator refuses the operation.  Missing
fields, duplicate declarations and malformed rational quantities fail earlier,
when the source is read.  The compiled record retains the original text hash
and the lines containing each operation and field.  Whitespace can change
that text hash without changing the compiled financial request.

\texttt{Annual\allowbreak{}Inventory\allowbreak{}Family/1} then gives each alternative its own name,
assumptions and ordered operations, all starting from the same declared annual
book.  Consider the earlier example with two recovered units costing 20 and
cash of 452.  The first alternative sells one unit for 15, assesses current tax
of 2, refreshes deferred measurement and carries the remaining unit into the
following year.  Its first-year profit is 3 and closing cash is 467.  Selling the
last unit for 15 in the next year gives that year's profit of 5 and closing cash
of 482, while the carried tax payable remains \(49/4\).  The two annual statement
vectors belong to this one history and are returned together with its name.

A different alternative may fail before the carry because the homogeneous-lot
finding is unknown, or may reach the second year and attempt to sell more units
than remain.  The family retains the precise valid prefix and the operations
that were not executed.  A report for a year that was never reached is
unavailable; it is not a statement with zero balances.  An explicit budget can
also leave members unvisited, while preserving their names and assumptions.
Even a zero member budget checks the common opening request.  Equal financial
results do not merge separately named alternatives.

Joint reports recompute the saved family before returning the complete archived
or current-prefix statements.  The reference checks cover inventory quantity,
current-tax histories, residual-stock carry, source locations, rejected premises
and retained alternatives.  These are authored, unweighted scenarios under the
selected proportional inventory model.  They establish neither probabilities
for those scenarios nor membership of an actual company's unknown history.
General inventory policies and independent accounting acceptance remain open.


\subsection*{Resume the annual book from its last committed operation}

A saved balance sheet is not enough to restart this simulation.  After the
first inventory sale, the remaining unit still carries its acquisition identity,
while each tax payable belongs to an assessment and retains any settlements.
The next operation must also be known: repeating the sale would change both
stock and cash even if the repeated journal balanced.  The optional
\texttt{Annual\allowbreak{}Inventory\allowbreak{}Store/1} therefore preserves the
complete annual request, the full period book and the ordered operation history.

The local store records the evaluated opening before attempting annual
operations.  Each later transaction records one attempted operation together
with its complete result and its position in the original request.  A successful
sale advances both the attempt position and the successful transaction history.
A refused sale advances the attempt position but leaves the last valid financial
state unchanged.  The refusal and all remaining operations stay visible, and
that immutable request stops.  Repairing its assumptions creates a different
request rather than rewriting the recorded history.

An explicit operation budget permits work to pause between transactions.
While operations remain unattempted, the stored result says that it is paused;
it does not turn a valid partial-year statement into a completed run.  A caller
can also supply the position it expected before a batch.  If that batch has
already committed, a retry returns the checked current result without applying
another sale.  If the caller expects a position that has not been reached, the
store rejects the request.

In the running example, the first sale leaves cash of 467, one recovered unit
costing 10, and the carried tax payable of \(41/4\).  Closing and reopening the
store retains all 26 accounts and their associated histories.  Continuing through
the declared tax assessment and period carry preserves the residual unit; its
sale in the following year gives cash of 482 and a payable of \(49/4\), exactly
as uninterrupted execution does.  The complete statements and refused prefixes,
not merely these selected balances, are compared in the reference checks.

Reopening checks the saved request, implementation identities, row order and
content hashes, then recomputes each committed prefix through the same financial
evaluator.  This catches altered balances even when a stored hash has also
been changed.  Replayed calculation has no external financial side effects.
A new operation uses the loaded full state directly; reopening deliberately
pays the cost of replay to validate that state.

The implementation uses one local SQLite writer and an atomic transaction for
each operation result.  In the process-exit checks, termination before commit
leaves the preceding result, while termination after commit retains exactly one
new result.  A fresh process can then reopen and continue.  These checks support
local restart behavior for this finite annual profile.  They do not establish
hardware power-loss durability, concurrent-writer correctness, efficient scaling
or authenticity against an attacker who controls the entire database.  Family
scheduling, migration and broader execution profiles remain separate work.



\section{Recorded statements before annual completion}
\label{sec:recorded-shared-statements}

The mixed return and lease sequence already determines useful financial
quantities. Its cash balance is 361, its lease liability has been paid, and its
recovered stock has been sold. Waiting for a complete annual tax calculation
should not make those recorded movements disappear. At the same time, a
balanced trial balance cannot establish that all required tax, valuation or
closing adjustments have been posted. The reporting question is therefore
precise: what do the completed events explain, and what remains unassessed?

\path{RecordedSharedBookStatements/1} answers that question for the synthetic
book. It replays the opening and requested event sequence, then constructs
financial views of the completed prefix. The supplied closing balance is never
an input. Each view identifies its last completed date, requested and completed
event counts, any failed event, and the declared fiscal period. Even when every
requested event succeeds, complete fiscal coverage and annual tax remain
unresolved. A short successful list does not prove that nothing else happened.

\subsection{A recorded balance sheet and equity movement}

The opening must start at the declared period boundary, with zero income
accounts and no prior event, cash-flow or settlement-tax-adjustment history in
the current execution. All 25 accounts receive exactly one classification.
The original return classes are extended by the right-of-use asset, lease
liability and lease gains; those additions are not hidden in a residual.
Owner-capital accounts must remain unchanged for these selected methods.

At the end of the five-event sequence, recorded assets total 701: cash 361,
inventory 60, property, plant and equipment 100, investments 80, deferred-tax
assets 10 and right-of-use assets 90. Other asset accounts are zero. Recorded
liabilities total 104, comprising debt 100 and deferred-tax liabilities 4.
The refund and lease liabilities are zero after their respective settlements.
All amounts are USD in the supplied synthetic example.

Recorded profit is revenue and gains less the expenses actually posted. The
sale contributes revenue 30 and the lease reduction a gain of 20. Cost of sales
is 16 after the settlement credit and subsequent sale, selling expense is 8,
interest is 11 and depreciation is 10. Current and deferred tax expense accounts
contain no postings in this sequence. Thus recorded profit is
\[
 P_{\mathrm{recorded}}=(30+20)-(16+8+11+10)=5.
\]
The zero tax postings describe the ledger; they are not a calculation of the
tax charge required for the period. Tax assessment stays unresolved alongside
the recorded result.

Let $A$ and $L$ be the totals of recorded assets and liabilities, and let $E_0$
be opening equity after reversing the debit-positive sign of the equity
accounts. The unchanged owner accounts give $E_0=592$. The two independent
reconciliations then agree:
\[
 A-L=701-104=597,
 \qquad E_0+P_{\mathrm{recorded}}=592+5=597.
\]
The first derives equity from the closing balance sheet; the second explains
its change from the recorded income. Any disagreement refuses the report. No
balancing reserve or owner contribution is inferred to make the equations hold.

\subsection{Explain the same cash movement in two ways}

Gross cash origins preserve the distinction between operating activity and
financing. Operating receipts are 30 from the recovered-stock sale. Operating
payments are 40 of refunds, 8 of handling and 11 of interest, totaling 59.
Lease principal of 110 is a financing payment under the retained historical
interest-classification assumption. There are no investing cash movements.
The direct cash view therefore gives
\[
 \mathrm{CFO}=30-59=-29,\qquad
 \Delta\mathrm{cash}=-29-110=-139=361-500.
\]
The complete gross flow records remain attached to the result, so these totals
can be traced back to the individual events rather than to a net cash plug.

The indirect view starts with recorded profit and removes its noncash effects.
Depreciation of 10 is added back. The gain of 20 from reducing the lease is
subtracted because the reduction produced no cash receipt. Releasing the
recovery asset of 16 adds 16 in the working-capital bridge. Settling the refund
liability reduces that credit balance from 40 to zero and uses 40 of cash.
Inventory returns to its opening amount after the sale, and the other selected
working-capital balances do not change. Consequently,
\[
 \mathrm{CFO}
 =\underbrace{5}_{\text{recorded profit}}
  +\underbrace{10}_{\text{depreciation}}
  -\underbrace{20}_{\text{noncash gain}}
  +\underbrace{16}_{\text{recovery asset released}}
  -\underbrace{40}_{\text{refund liability settled}}
 =-29.
\]
The sign of the refund term follows from the accounting convention. The signed
liability moves from $-40$ to zero; subtracting its change gives $-40$ in the
cash bridge. The credit-negative \path{other_gains} balance similarly removes
the positive noncash gain from profit. The reporting code uses the account
changes and gross flows to compute both reconciliations. The displayed values
are the checked example, not constants used to construct a report.

The inherited bridge also retains recorded deferred-tax expense and any
already-expensed opening interest paid as separate terms. The selected
successor does not execute those additional methods. They remain zero in this
example, and the fresh-opening requirement excludes prior paid interest from
the opening. Full tax measurement requires its own supported method and facts.

\subsection{A failed event leaves a smaller report, not a completed year}

In the alternative opening with cash 130, the return settlement leaves 82 and
the later payment of 121 fails. The recorded views describe only the lease
reduction and return settlement that completed. They show profit 16 and
operating cash outflow 48, with the failed payment and one unexecuted event
still visible. Their last completed date is 1 February 2024. They do not show
the lease as paid or describe the attempted payment date as a completed close.

If the first event is unsupported, a valid fresh opening can still be shown
with zero completed events and its explicit refusal. An invalid opening
supplies no financial views. An empty valid event list supplies the opening
view, but proves neither complete period coverage nor zero annual activity.
These distinctions preserve information without manufacturing a closing state.

All results keep tax assessment, period completion and annual completion
unresolved. The original native annual-report function retains its separate
tax-close prerequisite and still refuses the incomplete-tax control. The new
result is a recorded-prefix report with \path{annual_close_ready} and
\path{complete_financial_statements} false. Its arithmetic can be exact even
when required accounting adjustments remain unknown.

Complete replay checks statements, classifications, execution, gross origins,
source identity and coverage together. A changed profit, dropped flow, omitted
account, forged completion flag or altered component memory fails verification.
The implementation and complete requests/results are retained in
\path{scripts/shared_book_statements_v1.py} and
\path{data/accounting-master/shared-book-statements-20261024}.
This supplies the recorded-statement connection for the finite joint book;
complete annual methods and actual-source admission remain further work.


\section{Authoring the complete mixed book}
\label{sec:shared-book-authoring}

The return and lease example can now be stated as a complete text program.
An author supplies the opening accounts, the return obligation and inventory
lots, the lease rights, the premises attached to each method, and the ordered
events. The program ends with an explicit request for the recorded statements
of Section~\ref{sec:recorded-shared-statements}. These declarations lower to
the same typed request used by that example. No opening balance, missing
premise or additional event is supplied by the parser.

\subsection{Declare the opening before the events}

The opening has twenty-five signed account balances. All must be declared,
including accounts whose balance is zero. The inventory and lease collections
also have explicit opening and closing markers, even when empty. This makes
an omitted collection an authoring error, rather than evidence that the company
has no inventory lots or lease rights. The source description remains the
stated synthetic premise; writing its name does not turn it into an issuer fact.

For example, the cash account and the two empty inventory delimiters in the
checked opening appear as follows. This is an extract from the complete program:
the other twenty-four accounts and all other opening fields are required.
\begin{verbatim}
account "cash" "500/1";
lots;
end_lots;
\end{verbatim}
An amount such as \texttt{500/1} denotes an exact rational number. The same
representation carries rates, costs and cash payments; JSON floating-point
numbers are rejected. Structured values retain the return component, tax
premises, previous event identifiers and cash-flow history without reducing
them to account totals. The lease collection records each right's carrying
amount, liability, rate and last measurement date.

After the complete opening, the event block declares lease reduction, return
settlement, recovered-stock sale, lease payment and depreciation in that
order. Each event names the finite return or lease module and supplies its
full native record. The report declaration names the recorded-statement
profile and the exact source-bundle, opening and event digests. A changed
amount therefore requires a correspondingly changed request identity.
The compiler verifies that relationship; it does not authenticate the
economic premise behind the amount.

Complete programs, including every opening field and event parameter, are
saved under the shared-book-authoring evidence directory in the repository.
The accompanying grammar and typed record define the same bounded language.
Imports, include files, general expressions and user-supplied executable
functions are absent from this subset.

\subsection{Compilation and execution answer different questions}

Compilation asks whether the author supplied a well-formed program for the
selected methods. It checks closed field sets, exact scalar types, account
names and completeness, required collection markers, supported methods and
content digests. Duplicate JSON keys, accounts or lot keys are errors.
A failed declaration identifies its line. Whole-source checks first reject
invalid UTF-8, ambiguous line separators and excessive size or nesting.

Execution asks whether the declared opening and events support the proposed
accounting movements. An unbalanced opening can be structurally well formed,
yet fail the executor's balance check. Unknown lease membership remains
unknown after compilation. An unsupported currency remains unsupported.
Likewise, a duplicate event identifier or a payment without sufficient cash
is judged by the original executor, with the same completed prefix and
failure explanation as the typed request. Moving these judgements into the
parser would change the meaning of a successful compilation.

This distinction also preserves partial knowledge. A missing recovery estimate
that is explicitly recorded as unknown does not become zero. A depreciation
finding recorded as null does not become an affirmative finding. A false tax
premise remains false even if the recorded balance sheet balances. The text
program states those premises; the selected accounting method determines
what can follow from them.

\subsection{Preserve the whole result across representations}

The comparison uses all nineteen saved requests from the recorded-statement
example. Sixteen belong to this finite grammar. For each, text and typed
execution reproduce the complete original result: event trace, stopped prefix,
component memory, gross cash origins, classifications, reconciliations and
unresolved obligations. Equality of cash or profit alone would be insufficient.
The other three requests are explicitly rejected at compilation: one names an
unsupported tax module, one carries a stale opening digest, and one attempts
to supply an extra closing-state field. They are retained in the comparison
as refusals, rather than omitted from the evidence.

Text round trips preserve the complete typed program. Changing whitespace
changes the text digest but leaves the program and report-request digests
unchanged. Changing the program's name changes its program digest while
leaving the financial request unchanged. Thus source spelling, program
identity and the requested calculation remain distinguishable.

The cash competition remains visible through authoring. In the checked
130-cash opening, reduction and return settlement complete, leaving 82.
The subsequent 121 lease payment fails. The report preserves the first two
events and does not manufacture funds to finish the sequence. Separate checks
retain a nonselected opening inventory lot through all five mixed events,
and allow an explicitly empty lease collection with its separately supplied
balanced opening.

The resulting language is an executable authoring route for the declared
mixed book. Its complete-result comparisons support that bounded conclusion.
They do not settle the remaining language design for other accounting methods
or establish ease of use for human authors. Recorded statements still carry
the unresolved tax assessment and fiscal-period completion described in the
preceding section.


\section{Tax information before tax calculation}
\label{sec:shared-tax-information}

The mixed book has recorded profit, but it still lacks a tax assessment.
The earlier return-only reference does contain a tax-close method. Whether
that method can be reused is a substantive question: the new book contains
lease rights and a modification gain that the old book did not represent.
The unchanged method refuses the mixed state because its closed state shape
does not admit those additional accounts and lease records. The two original
tax examples still reproduce their complete results. Thus the refusal identifies
an incompatible extension, rather than a regression in the original method.

\subsection{A new income category changes the tax starting point}

The original pretax function sums the original income accounts, excluding
current and deferred tax. If that function is applied diagnostically to the
mixed balances, it returns a loss of 15. The complete recorded report instead
includes the lease-modification gain of 20. Under the signed-account convention,
the additional \texttt{other\_gains} balance is $-20$, so
\[
  P_{\mathrm{old\ accounts}}=-15,
  \qquad
  P_{\mathrm{mixed}}=-15-(-20)=5.
\]
The difference is an omitted income category, not a rounding error. Replacing
the state validator while leaving the original income list unchanged would
therefore calculate the wrong pretax starting point for the stated mixed-book
target. It would not preserve the meaning of the old calculation.

Even the corrected starting point does not determine tax. The return example
has a signed reconciliation adjustment of $-24$. Mechanical reuse of that
adjustment and the original illustrative rate would give
\[
  \frac{-15-24}{4}=-\frac{39}{4},
  \qquad
  \frac{5-24}{4}=-\frac{19}{4}.
\]
These expressions diagnose the effect of the omitted gain; neither is an
assessed tax or a cash refund. The lease deduction timing, the tax treatment
of the modification, the remaining right-of-use tax bases and the treatment
of a loss have not been supplied. The original finite method refuses its
negative-current-tax route. Copying the rate or reversing every deferred-tax
balance would not supply the missing facts.

Dates create a separate incompatibility. The original close on 31 December
stores 1 January as an exclusive tax endpoint. The mixed-book invariant
currently requires that endpoint to be no later than the state date. Those
two meanings cannot be combined by changing only the endpoint field. Moreover,
the mixed example's execution window runs from 1 January 2024 to 1 January 2026:
731 days. That window contains the selected lease and return events, but does
not establish a single annual tax interval or complete fiscal activity.

\subsection{Keep missing declarations and rejected premises visible}

A separate information assessor now accepts the unchanged recorded-report
request, an explicitly stated tax interval and identified information
declarations. The interval distinguishes its start, its last included calendar
day and its exclusive endpoint. The first finite implementation requires the
start to equal the recorded opening and every requested event to lie inside
the interval. The last recorded event may precede the asserted coverage date;
that assertion remains a declaration, not evidence that all activity is known.

For this mixed example, the assessor identifies 61 required declarations. They
cover jurisdiction and regime, the interval basis and coverage findings, rate,
loss and cash-settlement treatment, deferred recognition and recovery, every
pretax account's treatment and signed reconciliation adjustment, balance-sheet
tax bases, and the tax bases and timing of each lease right. Zero balances
remain in the account population. Both lease rights A and B remain relevant:
the native state retains fully terminated B with zero asset and liability,
and its earlier modification still requires a stated treatment.

Each supplied value carries an assumption or source-reference description.
Each unknown carries its reason. A source reference is retained as supplied;
the assessor does not authenticate it. With no information supplied, all
61 declarations remain unresolved. Removing any one from the hypothetical
complete control again leaves an unresolved requirement. A false coverage or
recognition finding remains false and is reported as unsupported. Extra fields,
wrong scalar types, forged request identities and incompatible interval
boundaries are invalid.

The complete native report accompanies the assessment. Its stopped event
prefix, component memory, gross cash movements and unresolved annual coverage
are preserved. When a requested payment fails, supplying every information
declaration does not erase that failure: completion of the requested sequence
remains unresolved. A first-year interval can be stated separately when its
requested events fit that interval, without relabelling the two-year example.

\subsection{Supplied information is not an executable tax rule}

The positive result is deliberately called \texttt{information\_supplied}.
It means the finite declaration set is present and has the required types.
The hypothetical controls exercise that representation; their supplied zero
bases and adjustments are not inferred from the financial statements. A text
field describing a deduction can still fail to specify a valid calculation.
Likewise, individually typed amounts can be inconsistent when aggregated.

Accordingly, every result keeps calculation readiness, source and accounting
admission, tax execution and annual-close readiness false. No tax amount is
posted. The next calculation design must state how account adjustments,
right-level tax bases, deferred recognition and settlements relate, and must
check their joint consistency. This preserves useful progress: missing facts
and competing declarations can be represented and examined now, while the
remaining numerical and accounting questions remain explicit.

\subsection{Reconcile a declared current-tax calculation}

The next finite calculation makes one further assumption explicit: a single
nonnegative rate applies to nonnegative taxable profit before credits or any
other regime-specific adjustments. This is a hypothetical arithmetic model,
named \texttt{uniform\_rate\_no\_credits\_v1}. It does not select a country's
tax rules. The author must select that model affirmatively; a false selection
is unsupported and an unknown selection remains unresolved.

Each declared account adjustment is now supported by identified allocation
items. An item taken from the completed return history must retain its original
event identity and signed amount, and must appear exactly once. Additional
items have explicit assumption references. The items assigned to an income
account must sum to that account's declared adjustment. Their category assignment
is supplied by the author; the arithmetic check does not establish its legal
appropriateness. In particular, the native return adjustment is not added again
after it has already been included in those account totals.

There are also four cross-field basis checks. At both opening and closing,
the declared asset bases of all lease rights must sum to the aggregate
right-of-use tax basis, and the declared liability bases must sum to the
aggregate lease-liability basis. The zero-valued right B participates in those
checks. A mismatch rejects the conditional calculation while preserving the
complete recorded financial report and the inconsistent declarations for
inspection. Neither carrying amount nor termination status supplies a tax basis.

\begin{samepage}
Let $P$ be recorded pretax income over all seven income accounts, $a_i$ the
identified signed adjustment items and $r$ the explicitly supplied rate. The
conditional taxable-profit expression and gross current tax are
\[
 T=P+\sum_i a_i,
 \qquad
 C=rT\quad\text{when }T\geq0.
\]
\end{samepage}
Here $C$ is gross conditional current tax under the selected arithmetic rule.
It is not an incremental expense, accrued payable or cash payment. The
calculation does not subtract opening accruals, apply credits, recognize
deferred tax or change the recorded closing state.

For the mixed example, $P=5$ and the retained return adjustment is $-24$.
Without any additional declared item, $T=5-24=-19$. The result is
\texttt{loss\_route\_required}; no negative tax amount or refund is computed.
A text declaration requesting a refund cannot change that result because it
does not supply an executable loss rule. Two hypothetical controls illustrate
the nonnegative branch. An expressly supplied additional adjustment of 19 gives
$T=0$ and $C=0$. An additional adjustment of 40 gives $T=21$ and, at the
illustrative rate $1/4$, $C=21/4$. Those extra adjustments are test assumptions,
not inferred company facts or recommended tax treatments.

The comparison retains the entire financial report, including cash of 361 and
recorded profit of 5, in every accepted arithmetic result. Inconsistent allocation
or basis declarations retain that same report alongside their rejection. Tests
also change native adjustment amounts, duplicate or omit identities, alter the
model finding and tamper with computed results. These checks establish the
specified arithmetic and preservation properties. A loss-utilization rule,
tax settlement, deferred recognition and annual close remain separate work.

\subsection{Carry a loss quantity without inventing a tax asset}

The negative expression of $-19$ needs a rule before it can affect a later
tax calculation. The finite extension supplies a conditional loss ledger.
Each lot has an identity, an origin reference, an availability date, an explicit
expiry declaration, original and remaining quantities, cumulative usage and
expiry, and an eligibility finding. An unknown finding remains unresolved;
a false finding is unsupported. A generated loss therefore becomes a quantity
in this separate ledger only under an affirmative, traceable assumption.
It is not a cash refund or a recognized deferred-tax asset.

The selected illustrative policy evaluates all declared usage on the last
included calendar day of an interval. An exclusive expiry on that day makes
the lot unavailable for usage. Expiry at the following endpoint permits usage
on the preceding day, but the remainder expires before the next opening.
No expiry is itself an explicit assumption. These conventions specify this
finite model; another jurisdiction or an intraperiod method may require a
different rule. Neither FIFO nor a legal utilization percentage is assumed.
The author supplies the chosen lot allocations and an interval usage cap.

For each lot, original loss equals remaining, used and expired quantities.
Across a step, with $L_0$ and $L_1$ denoting opening and closing remaining
quantities, $G$ generation, $U$ usage and $E$ expiry,
\[
 L_0+G=U+E+L_1.
\]
All quantities on this line are nonnegative exact rationals. A generated loss
becomes available at the exclusive endpoint and cannot be used in its own
generation interval. A usage cannot exceed its lot's remaining quantity,
the declared interval cap or positive taxable profit. Zero and exhausted lots
remain identified, preventing the same identity from being introduced again.
Consecutive ledger entries must agree on both interval endpoints and their
opening and closing quantities.

The checked negative example generates 19 under explicit eligibility and
no-expiry premises. A separate hypothetical 2026 calculation supplies taxable
profit of 21. Explicit use of 19 leaves taxable profit of 2 and gross conditional
tax of $1/2$ at rate $1/4$. This connects the loss ledger across periods; it does
not establish financial-book carry. The 2026 financial opening is independently
declared for the test and is expressly not the carried closing state of the
2024--2026 example. Both complete conditional calculations remain attached.

Partial use makes the expiry rule consequential. Using 7 of 19 leaves 12 when
the expiry is later than the next endpoint. If the exclusive expiry equals that
endpoint, 7 may be used on the preceding day and the remaining 12 expires.
If expiry instead equals the usage day, the usage is rejected; with no usage,
all 19 expires. These outcomes follow the selected policy, rather than a
silent decision to use the oldest lot or to preserve an expired balance.

The operation proposes a complete next loss state and per-lot movements.
Applying it recomputes the proposal and compares the complete result, including
the original financial calculation. Changed amounts, origins, eligibility,
scope, cursors or even Boolean-to-integer substitutions are rejected. This is
a pure in-memory application: it creates no durable commit, tax journal,
deferred asset or annual authorization. Connecting the conditional tax amount
and loss ledger to an actual financial tax balance still requires separately
specified recognition, accrual and settlement methods.

\subsection{Accrue current tax before choosing its payment}

A conditional calculation supplies an amount, but recognition and payment
change different financial balances. The next finite profile makes those
changes explicit. It begins with the complete recorded return-and-lease
calculation, recomputes the conditional loss result and requires an affirmative
recognition premise. An unknown premise remains unresolved; a false premise
does not produce a financial state. The profile applies only when the native
period has no current-tax expense already recorded and its tax-close cursor
still equals the assessment start. It therefore does not guess how an opening
tax accrual or an earlier assessment should be allocated.

The successor adds a separate income-tax payable to the original 25 accounts.
That liability is not trade payable. If the recomputed gross conditional tax
after loss usage is $C$, the accrual debits current-tax expense and credits
income-tax payable by $C$. Cash does not change. A zero amount produces a zero
journal. Existing deferred-tax assets, deferred-tax liabilities and deferred-tax
expense remain exactly as recorded; the loss balance does not create a new
deferred asset.

Payment requires its own identity, date, positive amount and affirmative
instruction. It debits income-tax payable and credits cash, subject to both the
outstanding liability and available cash. The author supplies the amount;
neither full payment nor a refund is inferred. A shortage of cash blocks that
payment while leaving a valid accrual available. This illustrative profile
classifies its specified tax payments as operating cash movements.

\begin{samepage}
Let $S$ be cumulative payment against the newly accrued amount $C$, with
$0\leq S\leq C$. Relative to the original recorded statements, the changes are
\[
 \Delta\text{profit}=-C,\qquad
 \Delta\text{tax payable}=C-S,\qquad
 \Delta\text{cash}=-S.
\]
\end{samepage}
Assets decrease by $S$, liabilities increase by $C-S$ and equity decreases by
$C$, preserving the balance-sheet identity. In the indirect operating cash
bridge, the lower profit and the increase in tax payable contribute
$-C+(C-S)=-S$. This equals the separately recorded cash payment. Paying the
liability does not charge the same expense again.

For the separate hypothetical 2026 calculation, loss usage leaves conditional
tax of $1/2$. Accrual records expense and payable of $1/2$ while cash stays at
500. An explicit payment of $1/4$ leaves cash of $1999/4$ and payable of $1/4$;
a second payment of $1/4$ clears the payable and leaves cash of $999/2$.
The recorded deferred-tax asset of 10 and liability of 4 are unchanged.
The original mixed example instead generates the loss quantity 19 and a zero
tax journal. These are consequences of the declared synthetic inputs, not
assessments of either a company or a jurisdiction.

The financial date and assessment coverage have separate meanings. Accrual is
dated on the last included day; coverage ends at the following exclusive
endpoint. The successor records that endpoint separately and leaves the native
tax-close cursor unchanged. Payments cannot precede the last tax operation or
fall outside the original financial period. An end-of-period accrual therefore
cannot silently accept a payment in the next year: financial-period carry must
first be specified.

The original native schema refuses the additional account and the exclusive
close cursor, so this is a distinct financial profile. Its validation replays
the retained calculation and every tax payment, then compares the complete
state. Changed balances, dates, loss results, payment histories and proposal
flags are rejected. Repeating the same accrual request produces the same pure
proposal; it does not establish globally unique durable posting. Later ordinary
return or lease events, financial-year carry, actual accounting admission and
annual completeness remain separate requirements.

\subsection{Continue ordinary events on the tax-adjusted book}

The additional tax payable must survive when ordinary business resumes. A
later return refund or lease payment has to use the cash left after tax
payments, and its journal must leave the tax balances intact unless that
method explicitly changes them. The continuation profile therefore retains
all 26 accounts, the return and lease memories, the assessment result and one
ordered event history. It extends the typed state and memory schemas directly;
it does not remove the tax payable to fit the older book.

The selected return methods instantiate the existing monetary definitions over
the complete account registry. Their memory operations still settle the return,
create the recovered lot or dispose of that lot, while preserving the additional
balances. The selected lease methods retain their original component equations
and map the resulting journal into the same registry. Tax expense and deferred
balances are unchanged by these ordinary methods. A separate, explicitly
instructed payment reduces income-tax payable and cash.

Assessment coverage constrains the next event. An ordinary event must occur on
or after the exclusive assessment endpoint and before the original financial
period ends. Otherwise it would change activity already used in the retained
tax calculation. A tax payment may occur on the accrual day, but cannot precede
an already completed operation. If assessment reaches the end of the declared
financial period, no ordinary continuation is possible without a separately
specified transition to the next opening.

\begin{samepage}
An earlier assessment provides a concrete same-period test. The author declares
the original two-year financial window, assesses the first interval through
31 December 2024 and places the later ordinary events on 1 January 2025.
The taxable expression of 2 and tax of $1/2$ are explicit hypothetical
assumptions. Paying that tax, settling the return, selling the recovered lot
and paying the lease use one cash balance:
\[
 500-\tfrac12-48+30-121=\tfrac{721}{2}.
\]
\end{samepage}
After the accompanying depreciation, recorded profit is $9/2$, assets are
$1401/2$, liabilities are 104 and equity is $1193/2$. Operating cash is
$-59/2$ and financing cash is $-110$. Direct cash, the indirect bridge,
the balance sheet and the equity movement reconcile exactly. The return's
later tax adjustment of $-24$ remains dated and identified for a subsequent
assessment; this continuation does not silently fold it into the earlier tax.

Payment order becomes consequential when cash is scarce. With opening cash
$485/4$, the lease payment of 121 alone leaves $1/4$. Paying tax of $1/2$
first leaves $483/4$, which cannot fund that lease payment. Reversing the order
allows the lease payment but leaves too little for tax settlement. Each request
stops at its first failed operation and retains the complete successful prefix.
It neither restores spent cash nor supplies financing automatically.

The result preserves the complete original report, loss state and conditional
tax calculation alongside the subsequent journals and gross cash movements.
Applying it recomputes the ordered prefix and compares the whole result.
Unknown findings remain unresolved, and repeated identities, changed balances
and backdated activity are refused. These checks establish selected same-period
execution under the declared assumptions. Further tax assessment, financial-year
carry and actual accounting admission remain separate work.

\subsection{Replace a provisional assessment without charging twice}

Later activity changes the question that the tax calculation must answer.
The continued book includes the return adjustment of $-24$, while its earlier
assessment covers only the preceding interval. A transition to the next
financial opening cannot make that unassessed activity disappear. The next
finite method therefore supplies an explicitly selected cumulative replacement:
a later calculation over the same original opening supersedes the earlier
provisional assessment. This is a hypothetical tax model. Independent tax
periods require a different method.

The replacement reads the actual continued 26-account book. It preserves the
cash already paid, retains the separate payable and includes every pretax
account and relevant lease right in the new declarations. Each native return
adjustment must appear exactly once with its event identity and amount.
Explicit additional assumptions must reconcile to the declared account totals,
and right-level tax bases must agree with their aggregate balances. Missing
declarations remain unresolved; incompatible scope, amounts or dates are refused.

The loss calculation returns to the original opening inventory. Using the
provisional closing ledger as a new opening could count a loss twice or make
it available inside the cumulative interval that generated it. The replacement
instead applies the declared usage, cap, eligibility and expiry rule once to
the original lots. The earlier result remains attached as superseded evidence.
For example, a provisional generated loss of 19 replaced by a cumulative loss
of 7 leaves 7, not 26. An original opening loss of 19, provisionally used by 1
and later used by 19 cumulatively, has total usage 19 rather than 20.

\begin{samepage}
Let $C_0$ be current-tax expense already recorded, $C_1$ the replacement total
after explicit loss usage, and $S$ actual tax payments. The additional expense
and resulting unpaid liability are
\[
 \Delta C=C_1-C_0,\qquad L_1=C_1-S.
\]
\end{samepage}
The method debits current-tax expense and credits tax payable by $\Delta C$.
A negative amount reverses those balances; it does not change cash. The
selected profile requires $L_1\geq0$. If earlier payments exceed the new total,
it preserves the calculation but refuses the financial proposal because a
tax receivable or refund needs a separately specified method.

In the paid mixed example, expense and payments are both initially $1/2$.
A declared replacement total of $21/4$ therefore adds only $19/4$ to expense
and payable. Cash stays at $721/2$. Pretax profit of 5 less total current-tax
expense of $21/4$ gives a recorded loss of $1/4$; assets of $1401/2$ less
liabilities of $435/4$ equal equity of $2367/4$. The operating cash result
remains $-59/2$. These amounts follow the explicit cumulative assumptions.

If no tax has been paid, replacing expense of $1/2$ with a total of $1/4$
reverses expense and payable by $1/4$ without moving cash. If $1/2$ was already
paid, the same replacement would leave a receivable of $1/4$ and is refused by
this profile. Both outcomes retain the complete continued financial history
and the replacement calculation for inspection.

The new coverage endpoint is later than the provisional endpoint and no later
than the original financial period boundary. The recorded events must fit
within it. Validation recomputes the continued book, the cumulative calculation,
the replacement losses and the incremental journal; applying a proposal compares
the whole result. Deferred-tax balances remain unchanged. Further assessment
chains, deferred-tax recognition, financial-period carry, source admission and
annual completeness still require their own methods and evidence.


\paragraph{Begin another period with recorded balances and unresolved findings.}
A conditional simulation can need another period before all the preceding
period's accounting questions have been resolved. For example, the scenario
above has cash, an unpaid current-tax liability, lease rights and a recorded
profit, but its deferred-tax assessment and the completeness of its historical
inputs remain open. Those recorded quantities are sufficient to construct a
provisional opening for selected subsequent transactions. They are insufficient
to declare that the preceding annual financial statements are complete.

The explicit model \texttt{provisional\_recorded\_balance\_carry\_v1} makes
that distinction part of execution. The request identifies the complete prior
reassessment, a unique carry event and the next period's start and exclusive
end. The prior current-tax assessment and loss history must reach the preceding
financial endpoint, and the next period must start at that endpoint. The prior
closing date must be its last included day. No calendar year or assessment of
missing transactions is inferred from those dates. Separate findings describe
event coverage, deferred-tax assessment, source admission and accounting
admission. Unknown and negative findings remain visible in the resulting state;
affirmative caller assumptions also leave annual readiness unestablished.

The transfer follows the book's debit-positive convention. Let $P$ denote the
recorded profit, $R$ the signed retained-earnings balance, and $b_k$ each
signed account balance immediately before the boundary. Let $\mathcal I$
contain the temporary income and expense accounts. The new opening satisfies
\[
 P=-\sum_{k\in\mathcal I}b_k,\qquad
 R^{\mathrm{open}}=R-P,\qquad
 b_k^{\mathrm{open}}=0\quad(k\in\mathcal I).
\]
Every other account keeps its recorded amount. A profit therefore increases the
credit balance in retained earnings; a loss reduces it. The transfer journal
sums to zero without a balancing plug. In the mixed scenario, recorded profit
is $-1/4$ and retained earnings is $-12$, so the opening signed balance is
$-12-(-1/4)=-47/4$. Cash stays $721/2$ and the unpaid tax liability stays
$19/4$. Positive-profit and zero-profit cases separately check the same sign
convention.

The boundary changes period accumulators without discarding economic memory.
Return components, recovered lots including disposed identities, lease rights,
measurement dates, loss quantities and lifetime event identifiers survive
exactly. The prior result retains its cash flows and tax adjustments together
with hashes of the closing book and report. The next period begins with empty
cash-flow and tax-adjustment lists, so an old receipt or deduction cannot be
counted again. Opening deferred-tax allocations copy the recorded asset and
liability; that copying does not assess their recognition or measurement. The
new period's assessment cursor denotes an empty interval and its current tax is
explicitly unassessed.

Subsequent statements use the actual carried opening. Paying the opening tax
liability of $19/4$ reduces cash from $721/2$ to $1423/4$, eliminates that
liability and leaves current-period tax expense at zero. In the indirect cash
calculation, the reduction in the opening payable explains the operating cash
outflow. It must not be charged to profit again. If depreciation of 10 follows,
recorded profit is $-10$, assets are $2743/4$, liabilities are 104 and equity is
$2327/4$. Depreciation is added back in the indirect bridge, leaving operating
cash flow at $-19/4$. The balance sheet, equity movement and direct and indirect
cash calculations reconcile to the same opening and closing state.

A second example carries an unsettled return component instead. Settlement
pays 48, releases a refund liability of 40 and a recovery asset of 16, and
recognizes recovered inventory of 20. Its cost-of-sales credit of 4 and selling
expense of 8 reduce profit by 4. Selling that inventory for 30 then earns 10,
so the two transactions together earn 6 and reduce cash by 18. Starting from
cash of 500, the following-period cash is 482. The indirect calculation gives
$6+16-40=-18$. The new return tax adjustment of $-24$ is retained for a later
assessment; it is not absorbed into an asserted tax-complete closing balance.

Ordered execution retains the completed prefix if a later event fails. A
payment still needs both an unpaid liability and actual shared cash; a reused
identity or an event outside the next period is refused. Applying a result
recomputes the entire predecessor and subsequent prefix and compares the
complete result, including unresolved findings and archived history. The
checked implementation covers this provisional boundary and selected return,
lease and tax-payment operations. New-period tax assessment, repeated carry,
complete event coverage and deferred-tax recognition require their own methods
and evidence before the scenario can support a complete annual close.


\paragraph{Measure deferred balances from explicit recognition assumptions.}
Carrying a recorded deferred-tax asset and liability preserves the book, but a
later assessment must explain whether their amounts should change. The mixed
book records an asset of 10 and a liability of 4. Its current-tax declarations
also contain tax bases. Neither fact establishes the amount of deductions that
can be recognized, the applicable rate on recovery or settlement, or whether a
movement belongs in profit or loss. The conditional deferred-assessment method
therefore requires those decisions separately and binds them to the actual
reassessed book.

The distinction follows the historical source treatment. IAS 12 paragraphs 7--9
distinguish asset and liability tax bases; paragraphs 24 and 27--29 condition
asset recognition on available taxable profit and its relevant grouping and
timing. Paragraphs 47 and 51 connect measurement to the applicable rate and
manner of recovery or settlement. Paragraphs 58 and 61A distinguish profit-or-loss
movements from tax related to items outside profit or loss
\citep[IAS 12, paras. 7--9, 24, 27--29, 47, 51, 58 and 61A]{eu20231803}.
These retained historical provisions explain the method's boundaries. The
example's rates, bases and recovery capacity remain explicit hypothetical
premises, rather than findings about an actual company's tax treatment.

The selected model uses one declared entity, tax authority and recovery group,
and keeps gross deferred assets and liabilities separate. Every nondeferred
asset and liability account is represented. Inventory is divided between other
inventory and retained lots; lease assets and liabilities are divided by right;
each remaining tax-loss lot is included separately. A component and its aggregate
cannot both contribute the same difference. Item bases must reconcile with the
complete prior declarations, while recovered-lot bases must also agree with
component memory. A finding about the absence of unrepresented tax-only items
and credits makes this finite scope explicit.

Let $C_i$ be the positive magnitude of an item's carrying amount, $B_i$ its
declared tax base and $r_i$ its declared nonnegative rate. An asset has signed
taxable difference $C_i-B_i$; a liability has $B_i-C_i$. A positive difference
is taxable and a negative one is deductible. For a loss lot, the supplied model
uses the negative of its remaining quantity, without consuming that quantity.
Write $d_i$ for the nonnegative deductible difference and $t_i$ for the
nonnegative taxable difference. The caller supplies a recognized deduction
$u_i$ and a common recovery capacity $K$, all in the same tax-base currency units:
\[
 0\leq u_i\leq d_i,\qquad \sum_i u_i\leq K,\qquad
 A_1=\sum_i r_i u_i,\qquad L_1=\sum_i r_i t_i.
\]
The sums include deductible temporary differences and loss lots together, so
one capacity cannot be reused for each item. The unrecognized amount $d_i-u_i$
remains visible. A deferred-tax liability does not automatically enlarge $K$:
the timing, entity and income restrictions needed to support recovery must be
supplied explicitly. The method does not produce that support by forecasting
future taxable profit.

Before changing balances, the request allocates the recorded asset $A_0$ and
liability $L_0$ completely across the same item population. Those allocations
must sum exactly to the book's recorded gross amounts. They explain inherited
balances rather than inventing a new assessment of the original opening.
For the supported profit-or-loss route, the change in deferred expense is
$(L_1-L_0)-(A_1-A_0)$. The journal debits the asset by $A_1-A_0$, credits the
liability by $L_1-L_0$, and records the balancing expense or benefit. Negative
changes reverse the corresponding balance. The request must support that
allocation; an equity, other-comprehensive-income or business-combination route
is refused by this particular implementation instead of being charged to profit.

In the checked example, property, plant and equipment has carrying amount 100
and tax base 60. At rate $1/4$, its taxable difference of 40 gives a liability
of 10. Debt has carrying amount 100 and tax base 60, so its difference is
deductible. Recognizing 32 of the available 40 gives an asset of 8 and leaves 8
unrecognized. Relative to the recorded asset of 10 and liability of 4, the asset
falls by 2 and the liability rises by 6; deferred expense is therefore 8.
Cash stays $721/2$, current tax and its unpaid liability are unchanged, and
recorded profit becomes $-33/4$. Assets of $1397/2$ equal liabilities of $459/4$
plus equity of $2335/4$. Adding the deferred expense back in the indirect cash
calculation leaves operating cash at $-59/2$.

Loss recognition uses the same capacity discipline. Recognizing deductions of
32 for the debt and 10 for a loss lot needs capacity of at least 42; capacity
of 32 is insufficient even if each item separately fits it. With capacity 42
and rate $1/4$, their combined deferred asset is $21/2$. The loss ledger still
contains the original remaining quantity of 19: recognizing a tax benefit does
not use the loss in the current-tax calculation. Matching bases, zero rates,
partial recognition and opposite asset/liability differences have separate
checked cases.

The proposal retains the entire prior result and changes only deferred balances,
deferred expense and its new event identity. Cash flows, current tax, loss
history, component state and assessment dates remain exact. Unknown necessary
findings stay unresolved, and an inconsistent population, allocation, base or
capacity produces no changed financial state. Pure apply compares the complete
recomputed proposal. Actual source and accounting admission, broader recovery
and allocation methods and annual completeness remain open. The preceding
frozen carry interface accepts a current-tax reassessment request; consuming
this newer result requires an explicit successor interface that preserves the
deferred journal and recognition schedule.


\paragraph{Carry the assessed book without renewing its recognition assumptions.}
The deferred journal changes the financial opening that a later period must
inherit. Replaying the older current-tax-only request would restore the asset
of 10 and liability of 4 and omit the deferred expense of 8. The resulting
retained earnings would also be wrong for the assessed scenario. The successor
carry therefore accepts the complete deferred-assessment request and recomputes
that proposal before transferring its actual closing balances.

The transfer uses the same debit-positive convention as the preceding carry.
Recorded profit is now $-33/4$, so retained earnings moves from $-12$ to
$-12-(-33/4)=-15/4$. The deferred asset of 8, deferred liability of 10, cash of
$721/2$ and current-tax payable of $19/4$ remain unchanged. Temporary income
accounts, including the deferred expense, begin at zero. Thus the next period
inherits the assessed economic amounts without charging the old expense again.

The complete predecessor remains attached. Its item-level differences,
recognized and unrecognized deductions, recorded allocations, support findings
and recovery capacity are retained with the original loss history. This archive
explains the opening balances; it does not renew the capacity or establish
new-period recognition. Opening deferred allocations copy the recorded 8 and
10, while both current and deferred tax in the next period are explicitly
unassessed. Unknown or negative carry findings remain visible independently of
the predecessor's conditional calculation.

If the next period pays the carried tax liability and records depreciation of
10, cash becomes $1423/4$ and recorded profit is $-10$. Assets of $2735/4$
equal liabilities of 110 plus equity of $2295/4$. Neither current nor deferred
tax expense is repeated. Operating cash is $-19/4$: depreciation is added back,
and the reduction in the opening tax payable explains the payment. The direct
cash calculation and the equity movement use this same carried opening.

The request retains the prior assessment and loss endpoints, requires a
contiguous next period and keeps lifetime event identities. Selected return,
lease and payment operations then act on the full new state. A failed operation
preserves the completed prefix. Applying a result compares the full recomputed
predecessor, opening, recognition archive and subsequent operations. The checked
scope remains provisional: a new current-tax assessment must distinguish
payments of the opening liability from payments of the new period's tax, and
new deferred recognition must reassess its own information. Repeated carry and
complete annual source and event coverage remain separate requirements.


\paragraph{Keep carried tax payments separate from the new assessment.}
The first current-tax assessment after carry has a different opening from the
same-period reassessment above. Current-period expense begins at zero, but the
book can still owe tax from the preceding period. Paying that opening liability
settles an old assessment; it does not reduce the new period's tax expense.
Calculating paid tax as current expense less the outstanding payable would
therefore be wrong for this opening.

The request binds the complete assessed-carry execution and begins its tax
interval at the new financial opening and loss cursor. It requires a completed
incoming execution and no previous assessment in the new period. Before that
first assessment, the selected payment kernel can consume only the carried
liability. Each payment can therefore be attributed to the prior assessment
from its exact event identity, date, amount and cash-flow record. A stopped
incoming execution is retained for inspection but is not silently assessed.

Let $P_0$ denote the opening tax payable, $S$ the completed payments against
that carried liability and $C$ the new conditional assessment. The two liability
records reconcile to the financial book as
\[
 P_{\mathrm{carried}}=P_0-S,\qquad
 P_{\mathrm{new}}=C,\qquad
 P_{\mathrm{total}}=P_0-S+C,\qquad 0\leq S\leq P_0.
\]
The first record retains the prior assessment's interval and identity with its
new-period settlements. The second records the new interval, assessment and
initially empty settlement history. The journal debits current-tax expense by
$C$ and credits the payable by the same amount. Cash, deferred balances and the
old settlements remain unchanged. Even a fully paid old liability keeps its
history rather than disappearing from the explanation.

For opening payable $19/4$ and new tax of 5, paying all $19/4$ leaves a total
payable of 5. Paying only 1 leaves $35/4$, and paying nothing leaves $39/4$.
The new expense is 5 in all three cases. In the fully paid example, after the
new-period depreciation of 10, recorded profit is $-15$, cash is $1423/4$,
assets are $2735/4$, liabilities are 115 and equity is $2275/4$. The increase
in the new payable cancels its expense in the indirect cash calculation; the
old payment remains the operating cash outflow of $19/4$.

Taxable profit is calculated from the new period's pretax accounts and explicit
adjustments. Each native return adjustment must appear exactly once, and the
complete tax-base declarations must reconcile at account and lease-right level.
Those declarations do not themselves assess deferred tax. In a return example,
new-period recorded profit of 6 and its native adjustment of $-24$ produce
taxable loss of 18. The method retains a new loss quantity of 18, available at
the interval's exclusive endpoint, without a refund or an automatic deferred
asset.

Loss use starts from the actual carried inventory. If a loss of 19 remains
available, explicit use against taxable profit of 20 leaves taxable profit of
1 and tax of $1/4$ at the declared rate. The history continues from the prior
cursor, with total use of 19. The selected close-day rule separately checks
expiry: a loss expiring before the usage day cannot be used, while a use before
expiry consumes only the declared quantity. Neither the previous deferred-tax
recognition nor its recovery capacity consumes or replenishes the loss inventory.

Applying the proposal compares the full predecessor, tax calculations, loss
movements, both liability records and financial statements. The checked scope
is a first assessment without credits, refunds or prepayments. Once both tax
vintages exist, subsequent payments need explicit allocation; the original
aggregate payment interface no longer supplies enough information. New-period
deferred assessment, repeated current assessments and later carry require their
own compatible methods. Source admission and complete annual event coverage
remain unresolved.


\subsubsection*{Name the tax liability settled by each payment}
Once carried tax and newly assessed tax are both payable, the payment amount
alone cannot identify which liability has been settled. A payment of 3 can reduce
either the carried payable of $19/4$ to $7/4$, leaving the new payable of 5
unchanged, or reduce the new payable to 2, leaving the carried payable at $19/4$.
Both choices reduce cash and total payable by 3. Their remaining tax obligations
have different histories, which the simulator must preserve.

The continuation therefore requires the caller to name each tax vintage and its
allocated amount. Every amount must be positive, no greater than the named
outstanding balance, and the allocations must sum to the payment. Actual shared
cash must cover that sum. An unknown authorization, a duplicated vintage, an
unknown liability or an excessive allocation prevents that operation. The
method supplies no first-in-first-out convention or proportional split.

For example, take the hypothetical first assessment through 30 June 2026 inside
the financial period ending on 31 December. Carried tax of $19/4$ is still
unpaid, new tax is 5, cash is $721/2$, and recorded profit is $-15$. Allocate 1
of a payment to carried tax and 2 to new tax. Cash becomes $715/2$; the two
payables become $15/4$ and 3, totaling $27/4$. Tax expense remains 5 and profit
remains $-15$. The payment is recorded in both affected settlement histories,
with their distinct amounts and the same dated cash movement.

An ordinary event can follow only from 1 July, the assessment's exclusive
endpoint. A further right-of-use depreciation charge of 2 reduces its asset from
80 to 78 and profit to $-17$, without changing either tax payable. Assets then
total $1367/2$, liabilities $467/4$, and equity including profit $2267/4$.
Operating cash flow remains $-3$. This depreciation is explicitly listed as
activity after the assessment; it has not received a new tax assessment.

The full book, prior loss inventory, deferred balances and unresolved findings
remain attached to the result. If a later operation fails, earlier completed
events and their settlement histories survive as a recorded prefix. Payments
can occur on the recorded assessment day, but an assessment at the financial
period's endpoint leaves no later date for ordinary events in that period.
Cumulative reassessment, new-period deferred recognition and another period
carry require their own compatible methods. The present calculation supplies
conditional bookkeeping evidence rather than annual accounting acceptance.


\subsubsection*{Replace the assessment while retaining the payments}
After further company activity, the new period's cumulative tax calculation may
change. Its recorded payments do not change with it. A replacement must therefore
start from the actual continued book and distinguish what was paid against the
carried liability from what was paid against the current period's assessment.
Using current expense less all tax payments would wrongly treat a payment of
last year's tax as a reduction in this year's expense.

Let $C_0$ be recorded current-period tax expense, $C_1$ the replacement cumulative
amount, $S$ the actual payments allocated to that current vintage, and $O$ the
remaining carried liability. The replacement records only the expense difference:
\[
  \Delta C=C_1-C_0,\qquad
  P_{\mathrm{current}}=C_1-S,\qquad
  P_{\mathrm{total}}=O+C_1-S.
\]
The journal debits current-tax expense by $\Delta C$ and credits payable by the
same signed amount. Cash and the carried liability remain unchanged. The current
vintage retains the identity referenced by earlier payments; its old assessment
and settlement record remain available alongside the replacement's provenance.

Continue the split-payment example above. Recorded current expense is 5,
current-vintage payments are 2, and carried payable is $15/4$. If the replacement
cumulative tax is 7, the expense increment is 2, current payable becomes 5 and
total payable becomes $35/4$. Cash stays at $715/2$ and recorded profit falls
from $-17$ to $-19$. Assets remain $1367/2$, while liabilities become $475/4$
and equity including profit becomes $2259/4$. If cumulative tax instead falls
to 3, the expense reduction is 2 and total payable becomes $19/4$.

A positive aggregate liability can conceal a different representation problem.
If cumulative current tax falls to 1 after payments of 2, the current vintage
would have a credit of 1 even though the carried liability leaves a positive
aggregate of $11/4$. The present positive-payable model refuses that replacement
and preserves the preceding book. A credit or refund needs a specified account
and treatment; it cannot be silently netted against the older obligation.

Losses are also recalculated cumulatively. Suppose the carried opening contains
an available loss of 19 and the first assessment uses all 19 against taxable
profit of 20. At a declared rate of one quarter, current tax is $1/4$. A
replacement using the same cumulative profit and deduction again yields $1/4$:
it restarts from the original opening inventory, rather than using the provisional
closing balance of zero or recording a second consumption of 19. Previous
provisional calculations remain in the history. If that loss expires before the
new assessment's usage date, the replacement must apply the stated expiry rule;
it cannot retain the earlier deduction automatically.

The cumulative interval must cover every completed event and cannot shorten the
previous assessment. A correction at the same endpoint is possible when no later
event lies outside it. Native return adjustments enter the replacement exactly
once; a worked return sequence has pretax income 6 and an adjustment of $-24$,
producing a new loss of 18 without an automatic deferred asset. Fresh deferred
recognition, another continuation or repeated replacement requires a compatible
method. Source admission, complete fiscal coverage and annual acceptance remain
unresolved.


\subsubsection*{Measure deferred tax from the newly assessed book}
The current-tax journal changes the book on which deferred tax will be measured.
In the cumulative example, payable was $27/4$ before the replacement and $35/4$
after it. The declarations used to prepare that current-tax calculation cannot
silently serve as measurements of the later book. The new deferred request
therefore names the exact recorded state and date, declares its jurisdiction and
regime, and supplies fresh bases for every scoped account and current lease
right. It also supplies the item-level bases and recognition decisions that
must reconcile to those totals. An unchanged base may be declared again with
its provenance; the method does not carry that premise forward automatically.

The selected hypothetical model retains the gross, single-recovery-pool,
profit-or-loss calculation described earlier. The complete current population
contains the ordinary asset and liability accounts, inventory subdivisions,
lease rights and remaining loss lots. Fresh measurement and recognition support
is required for every item. Recognized deductions across debt, losses and other
deductible items share the one explicitly supplied recovery capacity. The older
recognition schedule remains historical evidence. Its capacity cannot renew
itself, and a deferred liability cannot silently supply additional capacity.

Recorded allocations must also reconcile to the balances actually carried into
the period. In the worked book, those balances are a deferred tax asset of 8
and a deferred tax liability of 10. With property carrying amount 100 and a
fresh base of 80, the taxable difference is 20. With debt carrying amount 100
and base 60, the deductible difference is 40; suppose fresh support permits
recognition of only 20 against capacity 20. At a declared rate of one quarter,
the desired asset and liability are both 5. The journal reduces the asset by 3,
reduces the liability by 5, and records deferred-tax expense of $-2$.

After the preceding cumulative current-tax example, cash remains $715/2$,
current expense remains 7 and current-tax payable remains $35/4$. Recorded
profit rises from $-19$ to $-17$. Assets become $1361/2$, liabilities $455/4$
and equity including profit $2267/4$; operating cash flow is still $-3$.
All tax-vintage settlements, loss quantities, component states and earlier
findings remain unchanged. The result retains both the fresh recognition
premises and the previous period's schedule.

The capacity constraint matters when loss recognition is added. If the current
loss inventory has 9 remaining and the caller recognizes all 9 as well as the
20 debt deduction, capacity must cover 29. At one quarter, the combined desired
asset is $29/4$. Capacity of only 20 refuses that combination; the same amount
cannot support the debt deduction and then be used again for the loss. A loss
quantity by itself does not establish a deferred asset.

The method compares the entire proposed result when applied, including the
fresh measurements, item allocations, tax vintages and historical recognition
record. Unknown bases or support remain unresolved, and unsupported equity,
other-comprehensive-income and offset routes are refused. Another financial
boundary must preserve these histories while distinguishing old settlements
from the next period's cash movements. Conditional recognition here does not
establish annual completeness or source and accounting admission.


\subsubsection*{Cross the next boundary without replaying old payments}
The following-period book now contains current and deferred assessments, dated
payments against distinct tax liabilities, and the history of a replaced
assessment. Carrying this book into another period must preserve those histories
as well as its closing balances. The earlier carry caller cannot stand in for
that operation: it consumes a different predecessor and does not represent the
separate liability vintages.

The new boundary requires the recorded current and deferred assessments to reach
the financial endpoint, with the loss cursor at the same date and the closing
book dated the preceding day. The next financial interval must be contiguous.
This establishes the timing of a provisional carry, while unresolved source,
event-coverage and accounting findings remain attached to it. A caller's
assertion that those findings are true does not establish annual readiness.

In the fresh-deferred example, recorded profit is $-17$ and retained earnings
has a debit-positive balance of $-15/4$. Closing the income accounts therefore
sets retained earnings to $-15/4-(-17)=53/4$. All income accounts, including
current and deferred tax expense, start the new period at zero. Cash stays at
$715/2$, the deferred asset and liability stay at 5 each, and current-tax payable
stays at $35/4$. Assets, liabilities and total equity are unchanged. The boundary
journal balances, and the new opening has neither profit nor cash movement.

The two tax vintages retain their complete settlement and supersession records.
Separate new-period records carry outstanding amounts of $15/4$ and 5 under the
same liability identities. Each records the earlier vintage it came from, its
new financial period and an initially empty list of current-period settlements.
Earlier payments remain in the historical records; they do not enter the new
period's cash flows. A fully paid vintage is retained with a zero opening balance
rather than discarded.

The prior flows, tax adjustments, recognition calculation, measurement and
recovery capacity are archived with the closing book. Components, inventory lots,
lease-right dates, loss quantities and lifetime event identities remain intact.
Both current and deferred tax in the new period are unassessed; carrying a
recognition schedule does not renew its capacity or findings. This boundary
profile accepts no subsequent events. A compatible continuation must first
reconcile new payments against the new-period liability records while keeping
all earlier settlements historical. General annual execution and actual fiscal
completeness remain separate obligations.

\subsubsection*{Settle carried liabilities in the period that pays them}
The new opening can now support ordinary events and tax payments. Its two
liabilities still refer to earlier assessments, but the cash-flow statement
must record only payments made in the new financial period. Reusing the old
settlement list as current cash movement would count earlier payments again.
The continuation therefore keeps each historical vintage unchanged and records
new settlements against its separate current-period opening.

Suppose the company pays 3 on 1 January 2027 and explicitly allocates 1 to the
liability opening at $15/4$ and 2 to the liability opening at 5. Their remaining
amounts become $11/4$ and 3. Total payable falls from $35/4$ to $23/4$, while
cash falls from $715/2$ to $709/2$. The entry debits tax payable and credits cash
by 3; it creates no current-tax expense. The new period's operating cash flow
is $-3$. Historical payments remain available for inspection without entering
that sum.

Each payment names the stable liability identities and supplies positive exact
amounts whose sum equals the payment. Each allocation must fit that liability's
remaining current-period amount, and the total must fit available cash. The
caller also supplies affirmative payment authority. No oldest-first allocation,
proration, borrowing or reassessment is inferred. A second payment can exhaust
the remaining $11/4$ and 3, leaving both liability records at zero and cash at
$1395/4$. Keeping the zero records preserves which assessments were settled.

Ordinary events act on the same full financial state. If depreciation of 2
follows the first payment, the right-of-use asset falls from 78 to 76 and
recorded profit becomes $-2$. Assets are then $1351/2$, liabilities $443/4$ and
equity $2259/4$. The indirect operating cash calculation adds back depreciation
and includes the reduction in opening tax payable, giving $-3$, the same amount
as the dated cash payment. A retained inventory lot can also be sold through
the selected return method; its history and tax adjustments remain attached to
the resulting book.

Operations must have unused lifetime identities and chronological dates inside
the new financial interval. A rejected allocation, missing authority or cash
shortage stops that operation and preserves every preceding successful event,
together with the exact unprocessed sequence. Applying a proposed result
recomputes and compares the whole result, including its opening, financial
history, liability records and recognition archive. A nonportable request,
such as one containing a floating-point amount, is rejected before execution.

These events do not reassess current or deferred tax. The carried loss
quantities, recognition findings and historical measurements remain intact,
while both new-period assessments stay unassessed. A reusable annual design
must next explain how a fresh assessment joins all outstanding liabilities and
how another boundary preserves their complete histories.

\subsubsection*{Repeat the period with one explicit company book}
The next assessment may have several earlier liabilities to preserve. A
procedure that always separates one old liability from one current liability
cannot represent that situation. The common period book instead retains an
ordered collection of stable liability records, an immutable financial and
loss opening, current financial state, assessment histories and archived
periods. Each new period uses the same operations. The checked predecessor is
imported once, and a finite sequence of operations supplies the later events
and explicit assessment information.

A first current assessment adds one liability; a replacement changes that
liability's cumulative assessment while keeping its identity and payments.
Every earlier liability stays separate. Starting after the payment and
depreciation example, cash is $709/2$, the two liabilities are $11/4$ and 3,
and profit is $-2$. A hypothetical assessment of 4 adds a third liability,
raises total payable to $39/4$ and reduces profit to $-6$, without moving cash.
Paying 4 with allocations of 1, 1 and 2 then leaves cash at $701/2$ and the
liabilities at $7/4$, 2 and 2. If a fresh cumulative calculation replaces
the current assessment of 4 with 3, only the difference of $-1$ enters current
tax expense. The third liability becomes 1, total payable becomes $19/4$
and profit becomes $-5$. Replacing the assessment with less than the 2 already
paid requires a receivable or refund method; the selected no-credit model
refuses that operation.

Loss replacement starts from the immutable loss opening for the period. For
example, using 4 from an opening quantity of 9 leaves 5; replacing that
calculation with a use of 3 leaves 6, not 2. The earlier use remains in the
assessment history. Expiry is checked at the stated usage date, and a loss
identity cannot be reused after its earlier generation has been superseded
or carried into another period. Lease-right populations likewise include
identities observed in intermediate states, so a right that later disappears
is not silently removed from the period's tax-information requirements.

Fresh deferred measurement follows the current assessment and binds the actual
financial state. Its item bases, recorded balances, support and shared recovery
capacity must reconcile under the selected model. Another payment, native
event or current assessment makes that state-specific measurement stale. The
earlier calculation remains available, but the next boundary requires a new
measurement of the changed book. This dependency prevents an old recognition
capacity from renewing itself merely because the program continues.

At the financial endpoint, the current assessment and loss cursor must reach
the same exclusive date, the closing book must be dated the preceding day,
and the deferred measurement must bind that exact closing state. The
provisional carry closes actual profit into retained earnings and archives the
period's balances, reports, liabilities, payments, assessments and findings.
In the worked sequence, closing profit of $-5$ changes the debit-positive
retained balance from $53/4$ to $73/4$. Cash remains $701/2$. The next period
opens the three liabilities at $7/4$, 2 and 1, with no new settlements or cash
movements. Both tax assessments begin unassessed.

A new assessment of 2 in that period adds a fourth liability and raises total
payable to $27/4$. Profit is $-2$ and operating cash remains zero. The same
deferred-measurement and carry operations can then close that period, changing
retained earnings to $81/4$ while preserving all four liability identities.
These exact recurrences check the selected reference book and its refusal
behavior. Actual reporting completeness, admission of source evidence and
broader company or sector models still require their own specified methods
and evidence.

\subsubsection*{Author the same annual program in text and typed records}
The common period book also needs an authoring form that preserves its explicit
assumptions. A named annual program imports one checked root by its semantic
digest and lists operations in order. The import is supplied as a record in
the caller's local registry; it does not name a file to open or a service to
query. The compiler verifies that the imported record has the stated digest
and lowers the declarations to the same typed annual request.

After a first assessment creates the liability named
\texttt{assessment:annual-first}, a payment block can state an intended
allocation while preserving the fact that payment authority is unknown:
\begin{verbatim}
operation "allocated_tax_payment" "authored-payment";
date "2027-06-30";
amount "1/1";
authorization {"status":"unknown","reason":"Not supplied"};
payment_allocation {"vintage_id":"assessment:annual-first","amount":"1/1"};
end_operation;
\end{verbatim}
The enclosing program supplies the financial book and preceding assessment.
This declaration does not create the liability or authorize the payment. The
finite grammar uses one declaration per source line. Other
operation blocks expose native events, current-assessment information
and adjustments, loss use, deferred-measurement items, recovery capacity and
the next financial period as named fields or ordered entries. Empty
collections remain explicit in the typed request. No field supplies an
unstated accounting or behavioural assumption.

Structural compilation checks the closed operation shape, calendar dates,
exact rational values, declaration types and duplicate identities. It preserves
the distinction between known values with provenance and explicitly unknown
information. Whether an authorization is true, a payment fits available cash
or an assessment has sufficient information remains a financial evaluation
question. A program with an unknown tax rate therefore compiles and produces
the same unresolved assessment as its typed counterpart. A cash shortage
preserves the same successful prefix and unprocessed operations. Compilation
itself executes no financial method.

The compiler retains both the exact source hash and the semantic request hash,
with line locations for each operation and field. Adding blank lines changes
the source hash and locations but leaves the lowered request unchanged. The
canonical renderer reproduces the typed program, and both authoring paths
produce equal complete financial results for the checked scenarios. These
checks establish fidelity to the selected evaluator. They do not establish
that its supported operations exhaust the financial-statement language or
that a successfully parsed company program has achieved accounting admission.

\subsubsection*{Keep alternative annual histories and their statements together}
The same opening book can support several explicitly assumed sequences. A
company may have payment authority in one scenario and lack it in another;
the tax rate or recovery capacity may remain unknown. Comparing those
sequences requires keeping each company's balances, profit and cash movements
together. A cash balance from one scenario cannot be combined with the profit
or liabilities of another and presented as one company result.

An annual family names its members, states their assumptions and supplies one
common opening request and a consecutive reporting horizon. Each member lists
its own annual operations and planned period carries. The opening is bound by
its digest, and the planned periods must agree. Structural validation checks
these declarations without executing the financial methods. An unknown rate
can therefore remain in an otherwise well-formed member; the financial
evaluation determines where that member can proceed.

Each evaluated member uses the same annual evaluator as an individually
submitted program. Its complete financial history is retained. In the worked
sequence, the first period ends with cash $701/2$, recorded profit $-5$ and
operating cash flow $-7$. Its next period records profit $-2$, no cash movement
and total tax payable $27/4$. A fresh deferred measurement and another
provisional carry retain cash $701/2$ and change the debit-positive retained
balance to $81/4$. Both earlier periods remain available in the archive.

Other members need not reach those periods. If payment authority is false,
the first assessment succeeds but the payment is refused. The assessment,
its resulting balances and every remaining operation stay in that member's
result. Its later period has no report because execution never reached it.
An unknown rate stops before the first assessment; missing recovery capacity
or a stale measurement stops later. An explicit payment exceeding cash is
retained as a funding failure. None of these refusals removes the member or
prevents an independent member from continuing.

A joint period query returns the whole report for each member, bound to its
complete result. A period crossed by a successful provisional carry uses the
archived report. The current period uses its completed-prefix report and
states the actual date reached and the current and deferred assessment
statuses. Completing all listed operations does not turn that prefix into
an admitted annual statement. A period not reached is explicitly unavailable,
rather than a report of zeros. The query recomputes the family before trusting
saved results; changing cash and recomputing a content hash does not validate
an altered financial history.

The caller supplies an explicit member-evaluation budget and traversal order.
When that budget is exhausted, every unvisited member keeps its identity,
assumptions and request digest, with no invented financial result. Even a zero
member budget checks the common opening; that check is outside the count of
member evaluations. The count is not a guarantee about execution time or
memory. Increasing it makes a new request for recomputation, not a durable
continuation receipt.

These members are unweighted authored alternatives. Two names with different
stated assumptions remain distinct even when their financial results agree.
An empty authored selection does not establish that no feasible company
history exists. Evaluating every supplied member establishes coverage of that
list only. Reconstructing actual histories, proving an exhaustive feasible set
or attaching a probability law requires additional evidence and methods.


